\documentclass[11pt]{article}
% -------------
% Preambule pro DRILY (cvičné úlohy ve stylu statnic) - MFF UK, UI / Strojové učení
% Sazba: pdflatex (UTF-8, čeština). Samostatné PDF na úroveň (1/2/3).
% -------------
\usepackage[T1]{fontenc}
\usepackage[utf8]{inputenc}
\usepackage{lmodern}
\usepackage[czech]{babel}
\usepackage[a4paper,margin=2.3cm]{geometry}

\usepackage{amsmath,amssymb,mathtools}
\usepackage{textcomp}
\usepackage[nopatch=all]{microtype}
\usepackage{xcolor}
\usepackage{enumitem}
\usepackage{titlesec}
\usepackage[most]{tcolorbox}
\usepackage{booktabs}
\usepackage{listings}
\usepackage{hyperref}

% - seznamy -
\setlist{leftmargin=1.4em, itemsep=2pt, topsep=3pt, parsep=0pt}

% - barvy -
\definecolor{mffblue}{RGB}{27,54,93}
\definecolor{mffaccent}{RGB}{0,90,156}
\definecolor{ulohagray}{RGB}{70,70,75}
\definecolor{resgreen}{RGB}{30,110,60}
\definecolor{gapred}{RGB}{150,30,30}
\definecolor{calorange}{RGB}{180,110,0}
\definecolor{codebg}{RGB}{245,245,245}

\titleformat{\section}{\Large\bfseries\color{mffblue}}{\thesection}{0.6em}{}
\titleformat{\subsection}{\large\bfseries\color{mffaccent}}{\thesubsection}{0.6em}{}
\titlespacing*{\section}{0pt}{1.4em}{0.6em}
\setcounter{secnumdepth}{2}
\setcounter{tocdepth}{1}

\hypersetup{colorlinks=true, linkcolor=mffaccent, urlcolor=mffaccent,
  pdftitle={Cvičné úlohy ke státnicím - Informatika / UI / Strojové učení},
  pdfauthor={Viktor Číhal}, bookmarksnumbered=true}

\renewcommand{\arraystretch}{1.2}
\setlength{\parindent}{0pt}
\setlength{\parskip}{5pt plus 2pt minus 1pt}
\emergencystretch=3em
\hbadness=10000

% - kód (C#); komentáře držíme bez diakritiky kvuli listings+utf8 -
\lstdefinestyle{cs}{
  basicstyle=\ttfamily\footnotesize, backgroundcolor=\color{codebg},
  keywordstyle=\color{mffaccent}\bfseries, commentstyle=\color{resgreen!70!black}\itshape,
  stringstyle=\color{gapred}, showstringspaces=false, columns=fullflexible,
  breaklines=true, frame=single, framerule=0pt, rulecolor=\color{codebg},
  xleftmargin=6pt, xrightmargin=6pt, aboveskip=4pt, belowskip=4pt,
  literate={ě}{{\v e}}1 {š}{{\v s}}1 {č}{{\v c}}1 {ř}{{\v r}}1 {ž}{{\v z}}1
           {á}{{\'a}}1 {í}{{\'\i}}1 {é}{{\'e}}1 {ý}{{\'y}}1 {ů}{{\r u}}1 {ú}{{\'u}}1,
  language=[Sharp]C}
\lstset{style=cs}

% =====================================================================
% Prostředí pro ÚLOHU a ŘEŠENÍ ve stylu statnic.
% =====================================================================

% čítač úloh
\newcounter{uloha}
\renewcommand{\theuloha}{\arabic{uloha}}

% \uloha{okruh}{téma} -> nadpis úlohy ve stylu zadání statnic
% okruh: "společná matematika" / "společná informatika" / "specializace UI-SU"
\newcommand{\ulohahead}[2]{%
  \refstepcounter{uloha}%
  \par\medskip\noindent
  {\large\bfseries\color{ulohagray}Úloha \theuloha.}\quad
  {\small\colorbox{ulohagray!12}{\textbf{#1}}}\quad{\bfseries #2}\par
  \vspace{2pt}}

% zadání: lehce rámovaný blok jako na zkoušce
\newtcolorbox{zadani}{enhanced, breakable, sharp corners=south,
  colback=ulohagray!5, colframe=ulohagray, boxrule=0pt, leftrule=3pt,
  left=8pt, right=8pt, top=4pt, bottom=5pt, before skip=4pt, after skip=4pt}

% podotázky se značkou bodu (kolik bodu daná část typicky nese)
\newcommand{\bod}[1]{\tikz[baseline=-0.55ex]\node[circle,fill=mffaccent,text=white,
  inner sep=1.3pt,font=\bfseries\scriptsize]{#1};}
% bez tikz (rychlejší): kroužek přes textový symbol
\renewcommand{\bod}[1]{{\footnotesize\colorbox{mffaccent}{\textcolor{white}{\bfseries #1\,b}}}}

% řešení
\newtcolorbox{reseni}{enhanced, breakable, colback=resgreen!6, colframe=resgreen,
  boxrule=0.6pt, arc=1.5pt, left=7pt, right=7pt, top=3pt, bottom=4pt,
  before skip=4pt, after skip=8pt, fonttitle=\bfseries\small, coltitle=resgreen,
  title={Řešení}}

% kalibrační poznámka: jak tato úloha sytí podlahu (common>=5 / spec>=7)
\newcommand{\kalibrace}[1]{\par\smallskip\noindent
  {\footnotesize\textcolor{calorange}{\textbf{Kalibrace:}}~\textit{#1}}\par\smallskip}

% past / častá chyba
\newcommand{\past}[1]{\par\smallskip\noindent
  \textbf{\color{gapred}Pozor (častá chyba):}~#1\par}

% info box
\newtcolorbox{infobox}[1][]{enhanced, breakable, colback=mffaccent!5, colframe=mffaccent,
  boxrule=0.7pt, left=6pt, right=6pt, top=4pt, bottom=4pt, arc=2pt,
  before skip=6pt, after skip=6pt, fonttitle=\bfseries, #1}


\begin{document}

\begin{center}
  {\LARGE\bfseries\color{mffblue} Cvičné úlohy ke státnicím \\[2pt]
  Úroveň 1 \textnormal{(cíl: $>90\,\%$ pravděpodobnost složení)}}\\[4pt]
  {\small Informatika - Umělá inteligence, zaměření Strojové učení (UI-SU) \quad|\quad MFF UK}
\end{center}

\begin{infobox}[title={Jak je to zkalibrované (přečti první)}]
\footnotesize
Zkouška pro UI-SU = 8 otázek (2 spol.\ matematika, 2 spol.\ informatika, 4 specializace), každá 0-3 body.
Složení \textbf{není} hra o celkový součet, ale o \textbf{dvě nezávislé podlahy}:
\[
\text{projdeš} \iff \underbrace{(\text{společné}\ \ge 5/12)}_{\text{shovívavé: }2{+}1{+}1{+}1}\ \textbf{ a }\ \underbrace{(\text{specializace}\ \ge 7/12)}_{\text{přísné: }2{+}2{+}2{+}1}.
\]
Klíčový důsledek: ve specializaci musíš dostat \textbf{tři ze čtyř} otázek na $\ge 2$ body, kdežto společné okruhy ustojíš jednou jistotou plus definicemi. Proto úroveň~1 obsahuje:
\begin{itemize}\setlength{\itemsep}{1pt}
  \item plné zvládnutí nejčastějších témat (regrese, shlukování, evaluační metriky, rozhodovací stromy; lineární algebra, automaty);
  \item u \emph{každé} rodiny specializace (i ,,Základy UI'') aspoň \textbf{2-bodovou zálohu} = přesná definice + jeden kanonický postup;
  \item u společné informatiky ,,\textbf{nikdy nedostat 0}'' i na implementační otázce (binární soubor, semafor, OO v C\#);
  \item přesné definice/znění vět napříč okruhy (nejlevnější body).
\end{itemize}
\textbf{Jak procvičovat:} u každé úlohy zakryj \emph{Řešení}, odpověz nahlas a miř na ,,2 body'' (definice + jeden postup). Zelené pole \emph{Jádro na 2 body} u specializace je to, co musíš umět nazpaměť (cílí na to, abys nikdy nedostal 0-1 z pokryté otázky). Oranžová \emph{Kalibrace} říká, proč úloha sytí kterou podlahu.\\
\textbf{Proč $>90\,\%$ (a poctivě):} vázající je podlaha specializace; modelově $P(\text{spec}\ge7)\approx0{,}95$ (tři ze čtyř otázek na $\ge2$ body) a $P(\text{spol.}\ge5)\approx0{,}96$, tedy $P(\text{složím})\approx0{,}92$. Dvě věci, které tuto hodnotu zvedají a jsou už zapracované: pokrytí \emph{všech} rodin ,,Základy UI'' (i plánování, situační kalkulus, mechanism design, POMDP), aby vždy přítomný AI slot nepadl na nepokryté téma, a \emph{Jádro na 2 body} snižující šanci na skóre 0-1 u ML otázek.\\
\textbf{Aditivní struktura:} sešity se sčítají - úroveň 1 (tento) $\to{>}90\,\%$; $+$ úroveň 2 (výpočetní varianty, důkazové náčrty) $\to{>}95\,\%$; $+$ úroveň 3 (vzácný ,,ocas'' požadavku, zkoušky nanečisto) $\to\sim97,5\,\%$. Každý vyšší sešit obsahuje jen \emph{nové} úlohy, takže se nic neopakuje. Z 9 zkoušek dat a velkého oficiálního okruhu nelze poctivě garantovat $>99\,\%$; reálný strop pilného studenta je $\sim 97\text{-}98{,}5\,\%$.
\end{infobox}

\tableofcontents
\bigskip

\section{Společná matematika}
\ulohahead{společná matematika}{Vlastní čísla, diagonalizace a symetrické matice}
\begin{zadani}
Je dána reálná matice
\[
A=\begin{pmatrix} 2 & 1\\ 1 & 2 \end{pmatrix}.
\]
\begin{enumerate}[label=(\alph*)]
  \item \bod{1} Definujte vlastní číslo a vlastní vektor matice a charakteristický polynom. Uveďte vztah mezi vlastními čísly a determinantem (resp. regularitou) matice.
  \item \bod{2} Spočtěte vlastní čísla matice $A$ a ke každému najděte vlastní vektor.
  \item \bod{3} Rozhodněte, zda je $A$ diagonalizovatelná. Pokud ano, napište matice $P,D$ tak, že $A=PDP^{-1}$. Vysvětlete, proč je každá reálná symetrická matice (ortogonálně) diagonalizovatelná.
\end{enumerate}
\end{zadani}

\begin{reseni}
\textbf{(a) Definice.} Číslo $\lambda\in\mathbb{R}$ (obecně $\mathbb{C}$) je \emph{vlastní číslo} matice $A$, pokud existuje nenulový vektor $v$ s $Av=\lambda v$; takový $v$ je \emph{vlastní vektor}. Vlastní čísla jsou kořeny \emph{charakteristického polynomu} $p_A(\lambda)=\det(A-\lambda I)$. Platí $\det A=\prod_i \lambda_i$ (součin vlastních čísel s násobností), takže $A$ je regulární právě tehdy, když $0$ není vlastní číslo.

\textbf{(b) Výpočet.}
\[
\det(A-\lambda I)=\det\!\begin{pmatrix}2-\lambda & 1\\ 1 & 2-\lambda\end{pmatrix}=(2-\lambda)^2-1=\lambda^2-4\lambda+3=(\lambda-1)(\lambda-3).
\]
Vlastní čísla: $\lambda_1=3$, $\lambda_2=1$.
\begin{itemize}
  \item $\lambda_1=3$: $(A-3I)=\begin{pmatrix}-1&1\\1&-1\end{pmatrix}$, řešení $v_1=(1,1)^{\!\top}$.
  \item $\lambda_2=1$: $(A-I)=\begin{pmatrix}1&1\\1&1\end{pmatrix}$, řešení $v_2=(1,-1)^{\!\top}$.
\end{itemize}

\textbf{(c) Diagonalizace.} $A$ má dvě různá vlastní čísla, příslušné vlastní vektory jsou nezávislé, takže $A$ je diagonalizovatelná:
\[
P=\begin{pmatrix}1&1\\1&-1\end{pmatrix},\qquad D=\begin{pmatrix}3&0\\0&1\end{pmatrix},\qquad A=PDP^{-1}.
\]
$A$ je navíc symetrická. \emph{Spektrální věta}: každá reálná symetrická matice má reálná vlastní čísla a vlastní vektory příslušné různým vlastním číslům jsou na sebe kolmé, takže je lze znormovat na ortonormální bázi. Pak $A=QDQ^{\!\top}$ s ortogonální $Q$ ($Q^{-1}=Q^{\!\top}$). Zde $v_1\cdot v_2=1\cdot1+1\cdot(-1)=0$, takže $Q=\tfrac{1}{\sqrt2}\begin{psmallmatrix}1&1\\1&-1\end{psmallmatrix}$.
\end{reseni}

\kalibrace{Lineární algebra je nejčastější okruh společné matematiky (~65\,\% zkoušek). Část (a)+(b) = 2 body bezpečně sytí podlahu společných okruhů (cíl $\ge 5/12$). Definice z (a) je nejlevnější bod, nikdy ji nevynechej.}
\past{U diagonalizovatelnosti se neptáme jen ,,jdou spočítat vlastní čísla'', ale zda existuje báze z vlastních vektoru (tj. zda je algebraická = geometrická násobnost). Různá vlastní čísla to zaručí; u násobných je nutno ověřit dimenzi vlastního podprostoru.}

\ulohahead{společná matematika}{Soustavy rovnic, determinant, hodnost, skalární součin a ortogonalizace}
\begin{zadani}
\begin{enumerate}[label=(\alph*)]
  \item \bod{1} Definujte hodnost matice, regulární a inverzní matici a determinant. Uveďte vztah mezi regularitou, determinantem a vlastními čísly.
  \item \bod{2} Spočtěte determinant matice $B=\begin{psmallmatrix}1&2&0\\0&1&3\\1&0&1\end{psmallmatrix}$ a rozhodněte o její regularitě. Popište, jak Gaussova eliminace určí množinu řešení soustavy.
  \item \bod{3} Definujte skalární součin a kolmost. Gram-Schmidtovou ortogonalizací zpracujte vektory $v_1=(1,1,0)$, $v_2=(1,0,1)$ a výsledek znormujte.
\end{enumerate}
\end{zadani}

\begin{reseni}
\textbf{(a)} \emph{Hodnost} $\operatorname{rank}A$ je dimenze řádkového (= sloupcového) prostoru, tj.\ počet nenulových řádku v odstupňovaném tvaru. Matice $A\in\mathbb{R}^{n\times n}$ je \emph{regulární}, má-li hodnost $n$; pak existuje \emph{inverzní} $A^{-1}$ s $AA^{-1}=I$. \emph{Determinant} $\det A$ je multilineární antisymetrická funkce sloupcu s $\det I=1$; platí $\det(AB)=\det A\det B$, $\det A^{\!\top}=\det A$. Ekvivalentně: $A$ regulární $\iff \det A\neq 0 \iff 0$ není vlastní číslo $\iff \operatorname{rank}A=n$.

\textbf{(b)} Laplaceuv rozvoj podle prvního řádku:
\[
\det B = 1\cdot\det\!\begin{psmallmatrix}1&3\\0&1\end{psmallmatrix} -2\cdot\det\!\begin{psmallmatrix}0&3\\1&1\end{psmallmatrix} +0 = 1\cdot 1 -2\cdot(0-3) = 1+6 = 7 \neq 0,
\]
takže $B$ je regulární. \emph{Gaussova eliminace}: soustavu $Ax=b$ převedeme řádkovými úpravami na odstupňovaný tvar. Pak: pokud vznikne řádek $0=c$ s $c\neq0$, soustava nemá řešení; jinak je počet volných proměnných $=$ (počet neznámých) $-\operatorname{rank}A$ a řešení je jednoznačné právě tehdy, když $\operatorname{rank}A=$ počtu neznámých.

\textbf{(c)} \emph{Skalární součin} na reálném vektorovém prostoru $V$ je zobrazení $\langle\cdot,\cdot\rangle:V\times V\to\mathbb{R}$, které je \emph{symetrické} ($\langle x,y\rangle=\langle y,x\rangle$), \emph{bilineární} (lineární v každé složce) a \emph{pozitivně definitní} ($\langle x,x\rangle\ge0$, s rovností jen pro $x=0$). (V komplexním případě se požaduje hermitovskost/sesquilinearita.) Standardní příklad v $\mathbb{R}^n$ je $\langle x,y\rangle=\sum_i x_iy_i$. Indukovaná \emph{norma} $\lVert x\rVert=\sqrt{\langle x,x\rangle}$; vektory jsou \emph{kolmé}, je-li $\langle x,y\rangle=0$. Gram-Schmidt:
\[
u_1=v_1=(1,1,0),\quad u_2=v_2-\frac{\langle v_2,u_1\rangle}{\langle u_1,u_1\rangle}u_1=(1,0,1)-\tfrac12(1,1,0)=\big(\tfrac12,-\tfrac12,1\big).
\]
Kontrola: $\langle u_1,u_2\rangle=\tfrac12-\tfrac12+0=0$. Znormování: $e_1=\tfrac1{\sqrt2}(1,1,0)$, $\lVert u_2\rVert=\sqrt{\tfrac14+\tfrac14+1}=\sqrt{3/2}$, takže $e_2=\tfrac1{\sqrt6}(1,-1,2)$.
\end{reseni}

\kalibrace{Lineární algebra je nejčastější okruh společné matematiky; tato úloha jistí ,,výpočetní'' polovinu (soustavy/determinant/ortogonalizace) vedle úlohy o vlastních číslech. Definice z (a) = jistý 1 bod.}
\past{Determinant počítej rozvojem podle řádku/sloupce s nejvíce nulami. U Gram-Schmidta nezapomeň odečítat \emph{projekci} (s dělením $\langle u_1,u_1\rangle$), ne jen vektor.}

\ulohahead{společná matematika}{Limity, spojitost, derivace a Tayloruv polynom}
\begin{zadani}
\begin{enumerate}[label=(\alph*)]
  \item \bod{1} Definujte limitu funkce v bodě a spojitost funkce. Zformulujte větu o nabývání mezihodnot a větu o nabývání maxima na uzavřeném intervalu.
  \item \bod{2} Spočtěte $\displaystyle\lim_{x\to0}\frac{\sin x - x}{x^3}$. Pro $f(x)=x^3-3x$ najděte lokální extrémy a intervaly monotonie.
  \item \bod{3} Napište Tayloruv polynom (limitní formu) a pomocí něj zdůvodněte výsledek limity z (b). Určete konvexitu/konkavitu $f$.
\end{enumerate}
\end{zadani}

\begin{reseni}
\textbf{(a)} $\lim_{x\to a}f(x)=L$, pokud $\forall\varepsilon>0\,\exists\delta>0:\,0<|x-a|<\delta\Rightarrow|f(x)-L|<\varepsilon$. Funkce je \emph{spojitá} v $a$, je-li $\lim_{x\to a}f(x)=f(a)$. \emph{Věta o mezihodnotě (Bolzano):} je-li $f$ spojitá na $[a,b]$ a $y$ leží mezi $f(a),f(b)$, existuje $c\in[a,b]$ s $f(c)=y$. \emph{Weierstrass:} spojitá funkce na uzavřeném omezeném intervalu nabývá maxima i minima.

\textbf{(b)} Limita je typu $0/0$. l'Hospital třikrát, nebo lépe Taylor: $\sin x = x-\tfrac{x^3}{6}+o(x^3)$, takže
\[
\frac{\sin x-x}{x^3}=\frac{-x^3/6+o(x^3)}{x^3}\to-\tfrac16.
\]
Pro $f(x)=x^3-3x$: $f'(x)=3x^2-3=3(x-1)(x+1)$, nulové body $x=\pm1$. $f'>0$ na $(-\infty,-1)\cup(1,\infty)$ (rostoucí), $f'<0$ na $(-1,1)$ (klesající). V $x=-1$ lokální maximum, v $x=1$ lokální minimum.

\textbf{(c)} \emph{Tayloruv polynom} řádu $n$ v bodě $a$: $T_n(x)=\sum_{k=0}^n\frac{f^{(k)}(a)}{k!}(x-a)^k$, přičemž $f(x)=T_n(x)+o((x-a)^n)$ (limitní/Peanova forma). Rozvoj $\sin x$ kolem $0$ dává přímo čitatel $-x^3/6+o(x^3)$, odkud limita $-1/6$. Konvexita: $f''(x)=6x$, takže $f$ je konkávní na $(-\infty,0)$ a konvexní na $(0,\infty)$; v $x=0$ je inflexní bod.
\end{reseni}

\kalibrace{Analýza je druhý nejčastější okruh matematiky. Definice limity/spojitosti + jeden spočtený limit/extrém spolehlivě dají 2 body do podlahy společných okruhů.}
\past{l'Hospital lze použít jen na typy $\tfrac00$ nebo $\tfrac{\infty}{\infty}$ a po ověření, že limita podílu derivací existuje. Taylor je často rychlejší a méně chybový.}

\ulohahead{společná matematika}{Integrály, primitivní funkce a řady}
\begin{zadani}
\begin{enumerate}[label=(\alph*)]
  \item \bod{1} Definujte primitivní funkci a Riemannuv integrál a uveďte jejich vztah (Newtonova-Leibnizova formule). Definujte součet geometrické řady a uveďte vzorec pro $\sum_{n=0}^\infty ar^n$.
  \item \bod{2} Spočtěte $\displaystyle\int x e^{x}\,dx$ metodou per partes a určete $\displaystyle\int_0^1 x e^{x}\,dx$.
  \item \bod{3} Napište vzorec pro obsah pod grafem a pro objem rotačního tělesa. Vysvětlete odhad součtu řady integrálem (integrální kritérium).
\end{enumerate}
\end{zadani}

\begin{reseni}
\textbf{(a)} $F$ je \emph{primitivní funkce} k $f$ na intervalu, je-li $F'=f$. \emph{Riemannuv integrál} $\int_a^b f$ je společná limita horních a dolních součtu přes zjemňující se dělení (existuje, je-li $f$ např.\ spojitá). \emph{Newton-Leibniz:} je-li $F$ primitivní k spojité $f$, pak $\int_a^b f = F(b)-F(a)$. \emph{Geometrická řada:} $\sum_{n=0}^\infty ar^n=\frac{a}{1-r}$ pro $|r|<1$ (jinak diverguje); částečný součet $\sum_{n=0}^{N}ar^n=a\frac{1-r^{N+1}}{1-r}$.

\textbf{(b)} Per partes $u=x,\ dv=e^x dx$, tedy $du=dx,\ v=e^x$:
\[
\int x e^x\,dx = x e^x-\int e^x\,dx = (x-1)e^x + C.
\]
Určitě: $\int_0^1 x e^x\,dx=\big[(x-1)e^x\big]_0^1=(0)\cdot e-(-1)\cdot 1 = 1.$

\textbf{(c)} \emph{Obsah} pod nezápornou $f$ na $[a,b]$: $S=\int_a^b f(x)\,dx$. \emph{Objem} tělesa vzniklého rotací grafu kolem osy $x$: $V=\pi\int_a^b f(x)^2\,dx$. \emph{Integrální kritérium:} je-li $f$ kladná klesající, pak $\sum_{n=1}^\infty f(n)$ konverguje právě tehdy, když konverguje $\int_1^\infty f(x)\,dx$, a platí odhad $\int_1^{N+1}f \le \sum_{n=1}^{N} f(n) \le f(1)+\int_1^{N} f$.
\end{reseni}

\kalibrace{Integrál + geometrická řada jsou opakované matematické podotázky (geometrická řada byla i samostatná otázka). Vzorec per partes a součet geometrické řady umět zpaměti.}
\past{Per partes: vol $u$ tak, aby se derivováním zjednodušilo (polynom), $dv$ tak, aby šlo integrovat. Geometrická řada konverguje jen pro $|r|<1$ - vždy ověř.}

\ulohahead{společná matematika}{Teorie grafu: souvislost, stromy, barevnost, rovinnost, toky}
\begin{zadani}
\begin{enumerate}[label=(\alph*)]
  \item \bod{1} Definujte graf, stupeň vrcholu, souvislost a komponentu, strom (uveďte aspoň dvě ekvivalentní charakteristiky) a bipartitní graf.
  \item \bod{2} Definujte dobré obarvení a barevnost $\chi(G)$. Uveďte vztah barevnosti a klikovosti. Určete $\chi(C_5)$ (kružnice délky 5) a zdůvodněte.
  \item \bod{3} Zformulujte Eulerovu formuli pro rovinný graf a odvoďte horní mez počtu hran. Definujte síť, tok a řez a uveďte větu o max.\ toku a min.\ řezu.
\end{enumerate}
\end{zadani}

\begin{reseni}
\textbf{(a)} \emph{Graf} $G=(V,E)$, $E\subseteq\binom{V}{2}$. \emph{Stupeň} $\deg(v)$ = počet hran u $v$ (princip podání rukou: $\sum_v\deg v=2|E|$). $G$ je \emph{souvislý}, existuje-li cesta mezi každou dvojicí vrcholu; \emph{komponenta} = maximální souvislý podgraf. \emph{Strom} = souvislý graf bez kružnice; ekvivalentně: souvislý s $|E|=|V|-1$; nebo: mezi každými dvěma vrcholy existuje právě jedna cesta. \emph{Bipartitní} graf: $V$ lze rozdělit na dvě části tak, že každá hrana vede mezi nimi (ekvivalentně: neobsahuje lichou kružnici).

\textbf{(b)} \emph{Dobré obarvení} $k$ barvami: zobrazení $V\to\{1,\dots,k\}$ tak, že sousední vrcholy mají různou barvu. \emph{Barevnost} $\chi(G)$ = nejmenší takové $k$. Platí $\chi(G)\ge\omega(G)$ (klikovost; klika potřebuje tolik barev, kolik má vrcholu), ale rovnost obecně neplatí. $C_5$: jako lichá kružnice není bipartitní, takže $\chi>2$; třemi barvami obarvit lze, tedy $\chi(C_5)=3$.

\textbf{(c)} \emph{Eulerova formule:} pro souvislý rovinný graf s $V$ vrcholy, $E$ hranami a $F$ stěnami platí $V-E+F=2$. U jednoduchého grafu s $V\ge3$ je každá stěna ohraničena $\ge3$ hranami a každá hrana sousedí se 2 stěnami, takže $2E\ge3F$; dosazením $F=2-V+E$ dostaneme $E\le 3V-6$. (U grafu s mosty je třeba argument doplnit, závěr ale platí.) \emph{Síť} = orientovaný graf s kapacitami $c$, zdrojem $s$ a stokem $t$; \emph{tok} $f$ splňuje $0\le f\le c$ a Kirchhoffuv zákon (zachování v každém vnitřním vrcholu); \emph{řez} = rozdělení vrcholu na $S\ni s$ a $T\ni t$, jeho kapacita = součet kapacit hran z $S$ do $T$. \emph{Věta:} velikost maximálního toku se rovná kapacitě minimálního řezu (max-flow = min-cut).
\end{reseni}

\kalibrace{Grafy jsou stabilní okruh matematiky (barevnost, eulerovské/rovinné grafy, souvislost se v zadáních opakují). Definice + jedno odvození (např.\ $E\le3V-6$) = 2 body.}
\past{$\chi\ge\omega$ je jen nerovnost - existují grafy bez trojúhelníku s libovolně velkou barevností. Eulerova formule platí pro \emph{souvislé} rovinné grafy.}

\ulohahead{společná matematika}{Pravděpodobnost a statistika: Bayes, střední hodnota, rozptyl, CLT, testování}
\begin{zadani}
\begin{enumerate}[label=(\alph*)]
  \item \bod{1} Definujte pravděpodobnostní prostor, podmíněnou pravděpodobnost a nezávislost jevu. Zformulujte Bayesovu větu.
  \item \bod{2} Definujte střední hodnotu a rozptyl náhodné veličiny. Uveďte linearitu střední hodnoty a vzorec pro rozptyl součtu. Spočtěte $\mathbb{E}X$ a $\operatorname{var}X$ pro hod férovou kostkou.
  \item \bod{3} Zformulujte zákon velkých čísel a centrální limitní větu. Popište testování hypotéz: hladina významnosti, chyba 1.\ a 2.\ druhu.
\end{enumerate}
\end{zadani}

\begin{reseni}
\textbf{(a)} \emph{Pravděpodobnostní prostor} $(\Omega,\mathcal{F},P)$: množina elementárních jevu $\Omega$, $\sigma$-algebra jevu $\mathcal{F}$, míra $P$ s $P(\Omega)=1$. \emph{Podmíněná} $P(A\mid B)=\frac{P(A\cap B)}{P(B)}$ pro $P(B)>0$. Jevy \emph{nezávislé}, je-li $P(A\cap B)=P(A)P(B)$. \emph{Bayes:} $P(A\mid B)=\dfrac{P(B\mid A)P(A)}{P(B)}$, kde dle \emph{úplné pravděpodobnosti} $P(B)=\sum_i P(B\mid A_i)P(A_i)$ pro rozklad $\{A_i\}$ množiny $\Omega$ (po dvou disjunktní, pokrývající, $P(A_i)>0$).

\textbf{(b)} $\mathbb{E}X=\sum_x x\,P(X=x)$ (diskrétně), resp.\ $\mathbb{E}X=\int_{-\infty}^{\infty}x\,f_X(x)\,dx$ (spojitě). \emph{Linearita}: $\mathbb{E}(aX+bY)=a\mathbb{E}X+b\mathbb{E}Y$ (vždy, i pro závislé). $\operatorname{var}X=\mathbb{E}(X-\mathbb{E}X)^2=\mathbb{E}X^2-(\mathbb{E}X)^2$; $\operatorname{var}(aX)=a^2\operatorname{var}X$. Obecně $\operatorname{var}(X+Y)=\operatorname{var}X+\operatorname{var}Y+2\operatorname{cov}(X,Y)$, takže pro \emph{nezávislé} (kde $\operatorname{cov}=0$) je $\operatorname{var}(X+Y)=\operatorname{var}X+\operatorname{var}Y$. Kostka: $\mathbb{E}X=\tfrac{1+\dots+6}{6}=3{,}5$, $\mathbb{E}X^2=\tfrac{1+4+9+16+25+36}{6}=\tfrac{91}{6}$, $\operatorname{var}X=\tfrac{91}{6}-\tfrac{49}{4}=\tfrac{35}{12}\approx2{,}92$.

\textbf{(c)} \emph{ZVČ:} prumer $\bar X_n$ nezávislých stejně rozdělených veličin konverguje k $\mathbb{E}X$. \emph{CLT:} pro nezávislé stejně rozdělené se střední hodnotou $\mu$ a rozptylem $\sigma^2$ platí $\frac{\sqrt n(\bar X_n-\mu)}{\sigma}\xrightarrow{d}\mathcal{N}(0,1)$, tj.\ normovaný prumer má asymptoticky normální rozdělení. \emph{Testování:} nulová hypotéza $H_0$ vs.\ alternativa; \emph{hladina významnosti} $\alpha$ = max.\ přípustná pravděpodobnost \emph{chyby 1.\ druhu} (zamítnout platné $H_0$); \emph{chyba 2.\ druhu} $\beta$ = nezamítnout neplatné $H_0$; síla testu $=1-\beta$.
\end{reseni}

\kalibrace{Pravděpodobnost a teorie informace se v matematice objevují středně často; definice Bayes/EX/var + jeden výpočet je bezpečná záloha. Tyto pojmy navíc podpírají ML specializaci (regrese, naivní pravděpodobnostní úvahy).}
\past{Linearita střední hodnoty platí \emph{i pro závislé} veličiny; aditivita rozptylu jen pro \emph{nezávislé} (jinak $+2\operatorname{cov}$). Hladina významnosti řídí chybu 1.\ druhu, ne 2.}

\ulohahead{společná matematika}{Logika: normální tvary, sémantika, rezoluce}
\begin{zadani}
\begin{enumerate}[label=(\alph*)]
  \item \bod{1} Vysvětlete základní pojmy syntaxe výrokové a predikátové logiky (formule, otevřená/uzavřená formule). Definujte CNF, DNF a PNF.
  \item \bod{2} Definujte model, splnitelnost, tautologii a důsledek. Převeďte $(p\to q)$ na CNF a rozhodněte, zda je $(p\to q)\wedge p\to q$ tautologie.
  \item \bod{3} Popište rezoluci jako dokazovací metodu. Rezolucí ukažte, že množina klauzulí $\{p\vee q,\ \neg p,\ \neg q\}$ je nesplnitelná.
\end{enumerate}
\end{zadani}

\begin{reseni}
\textbf{(a)} \emph{Formule} se tvoří z prvovýroku (resp.\ atomických formulí s predikáty, termy a kvantifikátory) logickými spojkami. Ve výrokové logice není kvantifikace; v predikátové je formule \emph{otevřená}, má-li volnou proměnnou, a \emph{uzavřená} (sentence), nemá-li žádnou. \emph{CNF} = konjunkce disjunkcí literálu (klauzulí); \emph{DNF} = disjunkce konjunkcí literálu; \emph{PNF} (prenexní) = všechny kvantifikátory na začátku, za nimi otevřené jádro.

\textbf{(b)} \emph{Model} teorie/formule = ohodnocení (struktura), v němž je pravdivá. Formule je \emph{splnitelná}, má-li model; \emph{tautologie}, je-li pravdivá v každém ohodnocení; $\varphi$ je \emph{důsledek} teorie $T$ ($T\models\varphi$), platí-li ve všech modelech $T$. Převod: $p\to q\equiv\neg p\vee q$ (už je CNF, jediná klauzule). Výraz $((p\to q)\wedge p)\to q$ je \emph{modus ponens}; pravdivostní tabulkou (4 řádky) vyjde vždy $1$, je to tautologie.

\textbf{(c)} \emph{Rezoluce}: pracuje s klauzulemi (CNF); rezolventa dvojice klauzulí $C_1\vee\ell$ a $C_2\vee\neg\ell$ je $C_1\vee C_2$. Množina je nesplnitelná, lze-li z ní odvodit prázdnou klauzuli $\square$. Zde:
\[
\{p\vee q\},\ \{\neg p\}\ \Rightarrow\ \{q\};\qquad \{q\},\ \{\neg q\}\ \Rightarrow\ \square.
\]
Odvodili jsme prázdnou klauzuli, takže $\{p\vee q,\neg p,\neg q\}$ je nesplnitelná. (Doplňkově: věty o kompaktnosti a úplnosti zaručují, že rezoluce je úplná dokazovací metoda - detail je nad rámec úrovně 1.)
\end{reseni}

\kalibrace{Logika je vzácnější matematický okruh, ale definice (CNF/DNF, model, splnitelnost) jsou levné body a rezoluce/DPLL se prolínají i se specializací ,,Základy UI''. Tady stačí 2-bodová záloha.}
\past{Důsledek ($\models$) je sémantický pojem (přes modely), dokazatelnost ($\vdash$) syntaktický (přes odvození); věta o úplnosti je spojuje. Nezaměňuj splnitelnost (existuje model) s tautologií (všechny modely).}

\ulohahead{společná matematika}{Diskrétní matematika: relace, kombinatorika, inkluze-exkluze}
\begin{zadani}
\begin{enumerate}[label=(\alph*)]
  \item \bod{1} Definujte vlastnosti binární relace (reflexivita, symetrie, antisymetrie, tranzitivita), ekvivalenci a rozkladové třídy a částečné uspořádání (řetězec, antiřetězec).
  \item \bod{2} Kolik je prostých a kolik všech funkcí z $m$-prvkové do $n$-prvkové množiny? Zformulujte binomickou větu a princip inkluze a exkluze.
  \item \bod{3} Pomocí inkluze-exkluze odvoďte počet surjekcí z $m$-prvkové na $n$-prvkovou množinu a počet permutací bez pevného bodu (problém šatnářky).
\end{enumerate}
\end{zadani}

\begin{reseni}
\textbf{(a)} Relace $R\subseteq X\times X$ je \emph{reflexivní} ($\forall x: xRx$), \emph{symetrická} ($xRy\Rightarrow yRx$), \emph{antisymetrická} ($xRy\wedge yRx\Rightarrow x=y$), \emph{tranzitivní} ($xRy\wedge yRz\Rightarrow xRz$). \emph{Ekvivalence} = reflexivní + symetrická + tranzitivní; její \emph{rozkladové třídy} tvoří rozklad $X$. \emph{Částečné uspořádání} = reflexivní + antisymetrické + tranzitivní; \emph{řetězec} = po dvou porovnatelné prvky, \emph{antiřetězec} = po dvou neporovnatelné.

\textbf{(b)} Všech funkcí $X\to Y$ je $n^m$; \emph{prostých} (pro $m\le n$) je $n(n-1)\cdots(n-m+1)=\frac{n!}{(n-m)!}$. \emph{Binomická věta:} $(x+y)^n=\sum_{k=0}^n\binom{n}{k}x^k y^{n-k}$. \emph{Inkluze-exkluze:} $\big|\bigcup_{i=1}^n A_i\big|=\sum_i|A_i|-\sum_{i<j}|A_i\cap A_j|+\cdots+(-1)^{n+1}|A_1\cap\cdots\cap A_n|$.

\textbf{(c)} \emph{Surjekce} z $m$ na $n$: od všech $n^m$ funkcí odečteme ty, které vynechají aspoň jeden prvek cíle. Označíme-li $A_i$ funkce vynechávající $i$-tý prvek, dostaneme inkluzí-exkluzí
\[
\#\text{surjekcí}=\sum_{i=0}^{n}(-1)^i\binom{n}{i}(n-i)^m.
\]
\emph{Problém šatnářky} (derangement): permutace $\{1,\dots,n\}$ bez pevného bodu; $A_i$ = permutace s $\pi(i)=i$, $|A_{i_1}\cap\cdots\cap A_{i_k}|=(n-k)!$, takže
\[
D_n=\sum_{k=0}^n(-1)^k\binom{n}{k}(n-k)!=n!\sum_{k=0}^n\frac{(-1)^k}{k!}\approx \frac{n!}{e}.
\]
(Hallova věta o SRR a párování v bipartitním grafu je příbuzné téma - viz vyšší úroveň.)
\end{reseni}

\kalibrace{Diskrétní matematika dává levné definiční body (relace, ekvivalence) a jeden klasický I-E výpočet. Inkluze-exkluze je opakovaný typ.}
\past{Antisymetrie není ,,opak symetrie''. Počet surjekcí není $\binom{n}{?}$ trik - jde přes inkluzi-exkluzi; nezapomeň člen $i=0$ (= $n^m$).}


\section{Společná informatika}
\ulohahead{společná informatika}{Regulární jazyky: automaty, výrazy, uzávěry, neregularita}
\begin{zadani}
\begin{enumerate}[label=(\alph*)]
  \item \bod{1} Definujte deterministický (DFA) a nedeterministický (NFA) konečný automat, regulární gramatiku a regulární výraz. Zformulujte Kleeneho větu.
  \item \bod{2} Sestrojte DFA pro jazyk slov nad $\{a,b\}$ obsahujících podslovo $ab$. Uveďte uzávěrové vlastnosti regulárních jazyku (sjednocení, průnik, doplněk).
  \item \bod{3} Dokažte, že jazyk $L=\{a^n b^n\mid n\ge0\}$ není regulární.
\end{enumerate}
\end{zadani}

\begin{reseni}
\textbf{(a)} \emph{DFA} $=(Q,\Sigma,\delta,q_0,F)$ s $\delta:Q\times\Sigma\to Q$; přijímá $w$, skončí-li ve stavu z $F$. \emph{NFA} má $\delta:Q\times\Sigma\to 2^Q$ (a případně $\varepsilon$-přechody); přijímá, existuje-li přijímající běh. \emph{Regulární (pravá lineární) gramatika} je čtveřice $G=(N,\Sigma,P,S)$: neterminály $N$, terminály $\Sigma$, počáteční neterminál $S\in N$ a pravidla $P$ tvaru $A\to aB$ nebo $A\to a$ (případně $S\to\varepsilon$). \emph{Regulární výraz} se tvoří z $\emptyset,\varepsilon,$ symbolu operacemi $\cup,\cdot,{}^*$. \emph{Kleeneho věta:} třídy jazyku popsatelných DFA, NFA, regulární gramatikou a regulárním výrazem jsou totožné (= regulární jazyky).

\textbf{(b)} DFA se 3 stavy, $F=\{q_2\}$:
\[
\begin{array}{c|cc}
 & a & b\\\hline
q_0 & q_1 & q_0\\
q_1 & q_1 & q_2\\
q_2 & q_2 & q_2
\end{array}
\]
($q_0$ = zatím nic, $q_1$ = naposledy $a$, $q_2$ = už viděno $ab$). \emph{Uzávěry:} regulární jazyky jsou uzavřené na sjednocení a průnik (součinový automat na $Q_1\times Q_2$), na doplněk (u DFA prohodíme $F\leftrightarrow Q\setminus F$), dále na zřetězení a iteraci.

\textbf{(c)} Sporem přes \emph{pumping lemma}: nechť $L$ regulární s konstantou $p$. Vezmi $w=a^p b^p\in L$, $|w|\ge p$. Lemma dá rozklad $w=xyz$, $|xy|\le p$, $|y|\ge1$. Protože $|xy|\le p$, leží $y$ celé v úvodních $a$, tedy $y=a^k$, $k\ge1$. Pak $xy^2z=a^{p+k}b^p\notin L$ (více $a$ než $b$), což je spor. Tedy $L$ není regulární.
\end{reseni}

\kalibrace{Automaty a jazyky jsou nejčastější okruh společné informatiky (~78\,\% zkoušek, skoro vždy jako Q1). Konstrukce DFA + definice = spolehlivé 2 body; tato úloha je proto v jádru úrovně 1.}
\past{Pumping lemma dokazuje \emph{ne}regularitu (sporem); neslouží k důkazu regularity. Vždy zdůvodni, proč $y$ padne do bloku $a$ (díky $|xy|\le p$).}

\ulohahead{společná informatika}{Bezkontextové jazyky, Turinguv stroj a nerozhodnutelnost}
\begin{zadani}
\begin{enumerate}[label=(\alph*)]
  \item \bod{1} Definujte bezkontextovou gramatiku (CFG), jazyk jí generovaný a zásobníkový automat (PDA). Co je Chomského hierarchie (typy 0-3)?
  \item \bod{2} Sestrojte CFG pro $\{a^n b^n\mid n\ge0\}$ a popište princip převodu CFG na PDA.
  \item \bod{3} Definujte Turinguv stroj a rekurzivně spočetné jazyky. Vysvětlete existenci algoritmicky nerozhodnutelného problému (problém zastavení).
\end{enumerate}
\end{zadani}

\begin{reseni}
\textbf{(a)} \emph{CFG} $G=(N,\Sigma,P,S)$ s pravidly $A\to\alpha$, $A\in N$, $\alpha\in(N\cup\Sigma)^*$; \emph{generovaný jazyk} = slova nad $\Sigma$ odvoditelná z $S$. \emph{PDA} = konečný automat se zásobníkem; třída jazyku přijímaných PDA = bezkontextové jazyky (CFL). \emph{Chomského hierarchie:} typ 3 regulární $\subset$ typ 2 bezkontextové $\subset$ typ 1 kontextové $\subset$ typ 0 rekurzivně spočetné, s odpovídajícími stroji (KA, PDA, lineárně omezený automat, Turinguv stroj).

\textbf{(b)} CFG: $S\to aSb \mid \varepsilon$. \emph{Převod CFG$\to$PDA} (přijímání prázdným zásobníkem): na zásobník dej $S$; v každém kroku buď nahraď nejvyšší neterminál pravou stranou pravidla ($\varepsilon$-přechod), nebo paruj nejvyšší terminál se vstupním symbolem (a oba odeber). Slovo je přijato, vyprázdní-li se zásobník při dočtení vstupu. PDA tak simuluje levou derivaci.

\textbf{(c)} \emph{Turinguv stroj} = konečný řídicí automat s nekonečnou páskou, hlavou (čte/zapisuje/posouvá), přechodovou funkcí; přijímá vstup, dostane-li se do přijímajícího stavu. Jazyk je \emph{rekurzivně spočetný} (typ 0), přijímá-li ho nějaký TS (může cyklit na slovech mimo jazyk); \emph{rekurzivní} (rozhodnutelný), pokud TS vždy zastaví. \emph{Problém zastavení} (zastaví TS $M$ na vstupu $w$?) je nerozhodnutelný: kdyby existoval rozhodovač $H$, sestrojíme $D(x)$, který se zacyklí, právě když $H$ řekne ,,$x(x)$ zastaví''; pak $D(D)$ vede ke sporu (diagonalizace).
\end{reseni}

\kalibrace{Bezkontextové jazyky a Chomského hierarchie jsou jádro Q1; Turinguv stroj + nerozhodnutelnost jsou vzácnější, ale legální tah z téhož okruhu - tady je držíme na definiční úrovni jako pojistku proti nule.}
\past{,,Bezkontextové'' $=$ PDA, ne konečný automat. Nerozhodnutelnost se dokazuje diagonalizací/redukcí, ne ,,je to těžké''.}

\ulohahead{společná informatika}{Složitost, třídy P a NP, NP-úplnost}
\begin{zadani}
\begin{enumerate}[label=(\alph*)]
  \item \bod{1} Definujte časovou složitost a asymptotickou notaci ($O,\Omega,\Theta$). Definujte třídy P a NP.
  \item \bod{2} Definujte polynomiální převoditelnost, NP-těžkost a NP-úplnost. Uveďte aspoň tři NP-úplné problémy.
  \item \bod{3} Jak se dokazuje, že je problém NP-úplný? Co tvrdí Cookova-Levinova věta a co znamená otevřená otázka $\mathrm{P}\overset{?}{=}\mathrm{NP}$?
\end{enumerate}
\end{zadani}

\begin{reseni}
\textbf{(a)} \emph{Časová složitost} = počet elementárních kroku jako funkce velikosti vstupu (nejhorší případ). $f=O(g)$: $\exists c,n_0:\,f(n)\le c\,g(n)$ pro $n\ge n_0$; $\Omega$ je dolní mez, $\Theta=O\cap\Omega$. \emph{P} = problémy řešitelné deterministicky v polynomiálním čase. \emph{NP} = problémy, jejichž ,,ano'' instance mají v polynomiálním čase ověřitelný certifikát (ekvivalentně řešitelné nedeterministicky v polynomiálním čase).

\textbf{(b)} \emph{Polynomiální převod} $A\le_p B$: funkce vyčíslitelná v polynomiálním čase, která instance $A$ převádí na instance $B$ se zachováním odpovědi. $B$ je \emph{NP-těžký}, pokud $A\le_p B$ pro každé $A\in\mathrm{NP}$; \emph{NP-úplný}, je-li navíc $B\in\mathrm{NP}$. Příklady NP-úplných: SAT, 3-SAT, klika, vrcholové pokrytí, barevnost grafu, Hamiltonovská kružnice, batoh (rozhodovací).

\textbf{(c)} NP-úplnost se dokazuje ve dvou krocích: (1) ukázat $B\in\mathrm{NP}$ (certifikát + ověření); (2) zvolit známý NP-úplný problém $A$ a sestrojit polynomiální redukci $A\le_p B$. \emph{Cookova-Levinova věta:} SAT je NP-úplný (první NP-úplný problém), což dává startovní bod pro všechny další redukce. $\mathrm{P}=\mathrm{NP}$? je otevřená otázka, zda lze vše ověřitelné v polynomiálním čase i v polynomiálním čase \emph{vyřešit}; obecně se věří, že $\mathrm{P}\neq\mathrm{NP}$.
\end{reseni}

\kalibrace{Algoritmy a složitost (P/NP) jsou opakovaný okruh; čistě definiční úloha s vysokou návratností bodu. Redukce z (b)/(c) přidá třetí bod.}
\past{NP \emph{není} ,,ne-polynomiální''. NP-těžkost nevyžaduje $B\in$ NP; NP-úplnost ano. Směr redukce: ze \emph{známého} těžkého na \emph{nový} problém.}

\ulohahead{společná informatika}{Rozděl a panuj, Master theorem, třídění}
\begin{zadani}
\begin{enumerate}[label=(\alph*)]
  \item \bod{1} Popište princip metody rozděl a panuj. Zformulujte Master theorem (kuchařkovou větu) a dolní odhad složitosti porovnávacího třídění.
  \item \bod{2} Popište Mergesort (pseudokód) a odvoďte jeho složitost z rekurentní rovnice. Uveďte princip a nejhorší/průměrný případ Quicksortu.
  \item \bod{3} Pomocí Master theorem určete asymptotiku pro $T(n)=2T(n/2)+n$ a $T(n)=3T(n/2)+n$. Co z toho plyne pro Karatsubovo násobení?
\end{enumerate}
\end{zadani}

\begin{reseni}
\textbf{(a)} \emph{Rozděl a panuj}: rozděl úlohu na podproblémy, vyřeš rekurzivně, slož výsledky. \emph{Master theorem:} pro $T(n)=aT(n/b)+f(n)$ ($a\ge1,b>1$) porovnej $f(n)$ s $n^{\log_b a}$:
\begin{itemize}\setlength\itemsep{0pt}
  \item $f=O(n^{\log_b a-\varepsilon})\Rightarrow T=\Theta(n^{\log_b a})$;
  \item $f=\Theta(n^{\log_b a})\Rightarrow T=\Theta(n^{\log_b a}\log n)$;
  \item $f=\Omega(n^{\log_b a+\varepsilon})$ a regularita $\Rightarrow T=\Theta(f(n))$.
\end{itemize}
\emph{Dolní odhad} porovnávacího třídění: $\Omega(n\log n)$ (rozhodovací strom má $n!$ listu, hloubka $\ge\log_2 n!=\Theta(n\log n)$).

\textbf{(b)} Mergesort: rozděl pole na poloviny, setřiď rekurzivně, slij ($\textsc{Merge}$) v $O(n)$. Rekurence $T(n)=2T(n/2)+O(n)$, dle Master (případ 2) $T(n)=\Theta(n\log n)$; stabilní, $O(n)$ paměť navíc. Quicksort: zvol pivota, rozděl na menší/větší, rekurzivně; průměrně $O(n\log n)$, nejhůře $O(n^2)$ (špatný pivot), in-place.

\textbf{(c)} $T(n)=2T(n/2)+n$: $a=b=2$, $n^{\log_b a}=n$, $f=\Theta(n)$, případ 2 $\Rightarrow \Theta(n\log n)$. $T(n)=3T(n/2)+n$: $\log_2 3\approx1{,}585$, $n^{\log_2 3}$ roste rychleji než $f=n$, případ 1 $\Rightarrow \Theta(n^{\log_2 3})$. \emph{Karatsuba} násobí $n$-ciferná čísla rekurzí $T(n)=3T(n/2)+O(n)$, tedy $\Theta(n^{\log_2 3})\approx n^{1{,}585}$ - lepší než školní $\Theta(n^2)$.
\end{reseni}

\kalibrace{Třídění a rozděl-a-panuj jsou opakovaný algoritmický okruh; Master theorem + Mergesort je vděčný typový výpočet na 2-3 body.}
\past{Master theorem porovnává $f(n)$ s $n^{\log_b a}$, ne s $n$. Quicksort má nejhorší případ $O(n^2)$ - nezaměňuj s průměrným.}

\ulohahead{společná informatika}{Grafové algoritmy: prohledávání, nejkratší cesty, kostry, toky}
\begin{zadani}
\begin{enumerate}[label=(\alph*)]
  \item \bod{1} Popište prohledávání do šířky (BFS) a do hloubky (DFS), jejich složitost a reprezentaci grafu.
  \item \bod{2} Popište topologické třídění DAG a algoritmy pro nejkratší cesty: Dijkstra a Bellman-Ford (kdy který, složitost).
  \item \bod{3} Popište minimální kostru a algoritmy Jarník (Prim) a Borůvka (řezové lemma). Uveďte princip hledání maximálního toku (Ford-Fulkerson).
\end{enumerate}
\end{zadani}

\begin{reseni}
\textbf{(a)} \emph{BFS}: z počátku prohledává po vrstvách pomocí fronty; dává nejkratší cesty v neohodnoceném grafu. \emph{DFS}: jde do hloubky pomocí zásobníku/rekurze; klasifikuje hrany (stromové, zpětné, dopředné, příčné), detekuje cykly. Obě $O(V+E)$ při reprezentaci \emph{seznamy sousedu}.

\textbf{(b)} \emph{Topologické třídění} (jen DAG): lineární uspořádání vrcholu tak, že každá hrana vede ,,vpřed''; spočítá se opakovaným odebíráním vrcholu se vstupním stupněm 0 nebo přes časy dokončení DFS, $O(V+E)$. \emph{Dijkstra}: nejkratší cesty z jednoho zdroje pro \emph{nezáporné} váhy, hladově s prioritní frontou, $O((V+E)\log V)$. \emph{Bellman-Ford}: zvládá \emph{záporné} hrany, $V-1$ relaxací všech hran, $O(VE)$, navíc detekuje záporný cyklus.

\textbf{(c)} \emph{Minimální kostra} (MST) = kostra s minimálním součtem vah. \emph{Jarník/Prim}: roste jeden strom, vždy přidá nejlevnější hranu ven, $O((V+E)\log V)$. \emph{Borůvka}: každá komponenta paralelně přidá svou nejlevnější hranu, $O(E\log V)$. Korektnost plyne z \emph{řezového lemmatu}: nejlevnější hrana přes libovolný řez patří do nějaké MST. \emph{Ford-Fulkerson}: opakovaně hledej zlepšující cestu ze $s$ do $t$ v reziduální síti a posílej po ní tok, dokud existuje. Pro celočíselné kapacity vždy skončí; když už zlepšující cesta neexistuje, je tok maximální a jeho velikost se rovná kapacitě minimálního řezu.
\end{reseni}

\kalibrace{Grafové algoritmy jsou opakovaný informatický okruh (BFS/DFS, topo, MST, toky). Definice + jeden algoritmus s složitostí = 2 body; tady jde o breadth, ne hloubku důkazu.}
\past{Dijkstra selže na záporných hranách - tam Bellman-Ford. Řezové lemma je klíč ke korektnosti MST; znej ho aspoň slovně.}

\ulohahead{společná informatika}{Datové struktury: vyhledávací stromy a haldy}
\begin{zadani}
\begin{enumerate}[label=(\alph*)]
  \item \bod{1} Definujte binární vyhledávací strom (BVS) a jeho operace (find/insert/delete). Co je AVL strom a jakou má výšku?
  \item \bod{2} Popište binární (min-)haldu v poli a operace \textsc{Insert} a \textsc{ExtractMin}. Jak funguje Heapsort a jaká je jeho složitost?
  \item \bod{3} Popište hledání následníka (successor) v BVS. Porovnejte složitost operací u nevyváženého a vyváženého stromu. Proč jde halda postavit v $O(n)$?
\end{enumerate}
\end{zadani}

\begin{reseni}
\textbf{(a)} \emph{BVS}: binární strom, kde pro každý uzel platí klíče vlevo $<$ klíč uzlu $<$ klíče vpravo. \textsc{Find}/\textsc{Insert} jdou shora dle porovnání; \textsc{Delete} listu/uzlu s jedním synem je přímé, uzel se dvěma syny se nahradí svým následníkem. Vše $O(h)$ (výška). \emph{AVL} je BVS udržující v každém uzlu rozdíl výšek podstromu $\le1$ (rotacemi), takže $h=O(\log n)$.

\textbf{(b)} \emph{Min-halda} = úplný binární strom v poli (uzel $i$ má syny $2i{+}1,2i{+}2$), kde rodič $\le$ synové. \textsc{Insert}: přidej na konec a ,,probublej nahoru'' ($O(\log n)$). \textsc{ExtractMin}: vrať kořen, dej poslední prvek na vrchol a ,,probublej dolu'' ($O(\log n)$). \emph{Heapsort}: postav haldu a $n$-krát vyber minimum/maximum; $O(n\log n)$, in-place.

\textbf{(c)} \emph{Successor} (nejbližší větší): má-li uzel pravý podstrom, je to jeho nejlevnější uzel; jinak jdi nahoru, dokud nejsi levým synem rodiče (ten je následník). U \emph{nevyváženého} BVS je $h$ až $O(n)$ (degenerace na seznam), takže operace $O(n)$; u \emph{vyváženého} (AVL) $O(\log n)$. Haldu lze postavit \emph{zdola nahoru} (\textsc{Heapify} od posledního vnitřního uzlu): součet prací $\sum_h \tfrac{n}{2^{h+1}}\cdot O(h)=O(n)$.
\end{reseni}

\kalibrace{Stromy a haldy jsou opakovaný okruh datových struktur (BVS, AVL, Heapsort se objevily v zadáních). Definice + operace s složitostí = 2 body.}
\past{Mazání uzlu se dvěma syny v BVS jde přes následníka, ne ledabyle. Build-heap je $O(n)$, nikoli $O(n\log n)$ - klasická chytací otázka.}

\ulohahead{společná informatika}{Reprezentace dat a čtení binárního souboru (C\#)}
\begin{zadani}
Binární soubor má tento formát: 4 bajty ASCII ,,\texttt{REC1}'' (magic), poté \texttt{uint32} \texttt{count} (little-endian), poté \texttt{count} záznamů, kde každý záznam je \texttt{int32 id} (little-endian) následovaný \texttt{float32 value} (IEEE\,754, little-endian).
\begin{enumerate}[label=(\alph*)]
  \item \bod{1} Vysvětlete pojmy \emph{little-endian} vs.\ \emph{big-endian} a \emph{zarovnání} (alignment). Jak je v paměti uloženo \texttt{int32} s hodnotou \texttt{0x01020304} v little-endian?
  \item \bod{2} Napište v C\# metodu, která soubor načte do pole struktur \texttt{(int Id, float Value)}.
  \item \bod{3} Co se stane na big-endian stroji a jak to ošetřit? Jak ověřit magic a obránit se proti poškozenému/zkrácenému souboru?
\end{enumerate}
\end{zadani}

\begin{reseni}
\textbf{(a)} \emph{Endianita} je pořadí bajtu vícebajtové hodnoty v paměti. \emph{Little-endian} ukládá nejméně významný bajt na nejnižší adresu. Hodnota \texttt{0x01020304} se tedy v paměti jeví jako bajty \texttt{04 03 02 01} (od nízké adresy). Big-endian by uložil \texttt{01 02 03 04}. \emph{Zarovnání} znamená, že hodnota typu o velikosti $s$ má typicky začínat na adrese dělitelné $s$ (např.\ \texttt{int32} na násobku 4); překladač proto vkládá výplň (padding) mezi položky struktury. V binárním \emph{formátu souboru} se ale spoléháme na přesné pořadí bajtu a žádné implicitní zarovnání nepředpokládáme.

\textbf{(b)} \texttt{BinaryReader} v .NET čte v little-endian, což sedí na zadání:
\begin{lstlisting}
public readonly record struct Rec(int Id, float Value);

public static Rec[] Load(string path)
{
    using var fs = File.OpenRead(path);
    using var br = new BinaryReader(fs);

    // magic "REC1"
    var magic = br.ReadBytes(4);
    if (magic.Length != 4 || magic[0] != (byte)'R' || magic[1] != (byte)'E'
        || magic[2] != (byte)'C' || magic[3] != (byte)'1')
        throw new InvalidDataException("spatny magic");

    uint count = br.ReadUInt32();          // little-endian
    var recs = new Rec[count];
    for (uint i = 0; i < count; i++)
    {
        int id = br.ReadInt32();           // 4 bajty LE
        float val = br.ReadSingle();       // 4 bajty IEEE754 LE
        recs[i] = new Rec(id, val);
    }
    return recs;
}
\end{lstlisting}
Ruční složení int32 z bajtu (ukázka principu): \texttt{id = b0 | b1<<8 | b2<<16 | b3<<24}.

\textbf{(c)} Pozor na endianitu CPU: \texttt{BinaryReader} i explicitní složení \texttt{b0 | b1<<8 | b2<<16 | b3<<24} dávají \emph{vždy} little-endian výsledek nezávisle na procesoru (to je správně a přenositelné). Co přenositelné \emph{není}, je reinterpretace syrových bajtu v nativním pořadí CPU (např.\ \texttt{BitConverter.ToInt32} nebo přetypování ukazatele) - ta by na big-endian stroji vrátila prohozené bajty. Robustní je proto skládat bajty s pevným pořadím podle \emph{formátu} (LE): v .NET \texttt{BinaryPrimitives.ReadInt32LittleEndian(span)} a \texttt{BinaryPrimitives.ReadSingleLittleEndian(span)}, případně podle \texttt{BitConverter.IsLittleEndian} bajty otočit. \emph{Obrana proti poškození}: ověř magic; po načtení \texttt{count} zkontroluj, že zbývající délka souboru je $=8\cdot\texttt{count}$ bajtu (jinak vyhoď výjimku) místo slepého čtení do EOF; \texttt{ReadBytes}/\texttt{ReadInt32} při zkrácení vyhodí \texttt{EndOfStreamException}, kterou je nutno ošetřit.
\end{reseni}

\kalibrace{Implementační otázky ze společné informatiky (binární soubor, alokátor, semafor, OO návrh) jsou hlavní příčina propadnutí ,,naoko připraveného'' studenta - na nule z této otázky se podlaha $\ge 5/12$ hroutí. Cíl: NIKDY nedostat 0. Část (a) je čistá definice (1 bod), kostra metody v (b) přidá druhý bod i bez doladění (c).}
\past{\texttt{BinaryReader} čte vždy little-endian bez ohledu na CPU - na zkoušce to klidně řekni, ale dodej, že pro přenositelnost a big-endian CPU je správně \texttt{BinaryPrimitives}. Nezapomeň na ošetření konce souboru.}

\ulohahead{společná informatika}{Architektura a OS: režimy, procesy, virtuální paměť}
\begin{zadani}
\begin{enumerate}[label=(\alph*)]
  \item \bod{1} Vysvětlete adresu a adresový prostor, privilegovaný vs neprivilegovaný režim procesoru a roli jádra OS. Co je proces, vlákno, kontext a přepínání kontextu? Uveďte stavy vlákna.
  \item \bod{2} Vysvětlete rozdíl virtuální vs fyzická adresa a komponenty stránkování (číslo stránky, číslo rámce, offset). Pro stránku 4\,KiB přeložte virtuální adresu \texttt{0x00003ABC}, je-li stránka 3 namapována na rámec 7.
  \item \bod{3} Jakou roli hrají virtuální adresové prostory v ochraně paměti? Jak procesor realizuje konstrukce vyššího jazyka (přiřazení, podmínka, cyklus) instrukcemi?
\end{enumerate}
\end{zadani}

\begin{reseni}
\textbf{(a)} \emph{Adresa} identifikuje místo v paměti; \emph{adresový prostor} = množina platných adres (každý proces má vlastní). V \emph{privilegovaném} (jádrovém) režimu smí procesor vše (I/O, správa paměti), v \emph{neprivilegovaném} (uživatelském) jen omezeně; přechod přes přerušení/syscall. \emph{Jádro OS} spravuje prostředky a izolaci. \emph{Proces} = běžící program s vlastním adresovým prostorem; \emph{vlákno} = tok provádění uvnitř procesu (sdílí paměť). \emph{Kontext} = registry + čítač instrukcí + stav; \emph{přepnutí kontextu} je uloží/obnoví. Stavy vlákna: \emph{běží} (running), \emph{připraveno} (ready), \emph{čeká/blokováno} (waiting); plánovač přepíná preemptivně nebo kooperativně.

\textbf{(b)} \emph{Virtuální} adresu používá program; \emph{fyzická} ukazuje do RAM; překlad zajišťuje MMU přes stránkovací tabulku. Adresa $=$ (číslo stránky $\Vert$ offset). Stránka 4\,KiB $\Rightarrow$ offset 12 bitu. \texttt{0x00003ABC}: offset $=$ \texttt{0xABC}, číslo stránky $=$ \texttt{0x00003ABC}$\gg12=$ \texttt{0x3} $=3$. Tabulka: stránka $3\to$ rámec $7$. Fyzická adresa $=(7\ll12)\,\vert\,\texttt{0xABC}=\texttt{0x7ABC}$.

\textbf{(c)} Každý proces má vlastní mapování, takže nevidí cizí paměť (\emph{ochrana/izolace}); přístup mimo mapování vyvolá výpadek (page fault). Konstrukce vyššího jazyka: \emph{přiřazení} = load/store mezi registrem a pamětí; \emph{podmínka} = porovnání (\texttt{cmp}) + podmíněný skok; \emph{cyklus} = podmíněný skok zpět; \emph{volání funkce} = uložení návratové adresy a rámce na zásobník (\texttt{call}/\texttt{ret}).
\end{reseni}

\kalibrace{Architektura a OS jsou druhý nejčastější informatický okruh (~60\,\%); virtuální paměť a překlad adres se v zadáních opakují jako konkrétní výpočet. Tabulka adres je vděčné 2-3 body.}
\past{Offset má tolik bitu, kolik je $\log_2$(velikost stránky); nepleť číslo stránky (virtuální) s číslem rámce (fyzické). Stavy vlákna nejsou jen ,,běží/neběží''.}

\ulohahead{společná informatika}{Synchronizace: kritická sekce, semafory, producent-konzument (C\#)}
\begin{zadani}
\begin{enumerate}[label=(\alph*)]
  \item \bod{1} Definujte časově závislou chybu (race condition), kritickou sekci a vzájemné vyloučení. Co je zámek a jaký je rozdíl mezi aktivním a pasivním čekáním?
  \item \bod{2} Co je semafor (operace \textsc{wait}/P a \textsc{signal}/V)? Napište v C\# řešení producent-konzument s omezeným bufferem pomocí semaforu.
  \item \bod{3} Uveďte čtyři podmínky uváznutí (deadlock). Porovnejte spinlock a mutex a vysvětlete atomickou operaci compare-and-swap (CAS).
\end{enumerate}
\end{zadani}

\begin{reseni}
\textbf{(a)} \emph{Race condition}: výsledek závisí na nedeterministickém prokládání vláken nad sdíleným stavem. \emph{Kritická sekce}: úsek kódu, který smí běžet nejvýše jedno vlákno najednou. \emph{Vzájemné vyloučení} (mutual exclusion) tuto vlastnost zajišťuje. \emph{Zámek} (lock) chrání kritickou sekci; \emph{aktivní čekání} (spin) cyklí a pálí CPU, \emph{pasivní} vlákno uspí a plánovač ho probudí (lepší při delším čekání).

\textbf{(b)} \emph{Semafor} = čítač s atomickými operacemi: \textsc{wait}/P sníží čítač (a zablokuje při 0), \textsc{signal}/V zvýší (a probudí čekající). Omezený buffer kapacity $N$:
\begin{lstlisting}
// empty: kolik je volnych mist, full: kolik je obsazenych
var empty = new SemaphoreSlim(N, N);
var full  = new SemaphoreSlim(0, N);
var mutex = new object();
var buffer = new Queue<int>();

void Producer(int item) {
    empty.Wait();                 // pockej na volne misto (P)
    lock (mutex) { buffer.Enqueue(item); }
    full.Release();               // pribyl prvek (V)
}

int Consumer() {
    full.Wait();                  // pockej na prvek (P)
    int item;
    lock (mutex) { item = buffer.Dequeue(); }
    empty.Release();              // uvolnil se slot (V)
    return item;
}
\end{lstlisting}

\textbf{(c)} \emph{Deadlock} nastane při současném splnění: (1) vzájemné vyloučení, (2) drž a čekej, (3) bez odebrání (no preemption), (4) kruhové čekání. \emph{Spinlock} aktivně čeká (vhodný pro velmi krátké sekce na více jádrech), \emph{mutex} uspí vlákno (vhodný pro delší). \emph{CAS}\,(\texttt{CompareAndSwap}(addr, expected, new)): atomicky porovná a případně zapíše; základ neblokujících (lock-free) struktur.
\end{reseni}

\kalibrace{Synchronizace (semafory, kritická sekce, race condition) je nejčastější OS podtéma a často chce \emph{kód}. Definice (a) = jistý bod, kostra semaforu (b) = druhý; tím se vyhneš nule na implementační otázce.}
\past{Pořadí semaforu u producent-konzument je klíčové: čekej na zdroj (\textsc{wait}) PŘED zámkem, jinak hrozí deadlock. \texttt{empty} a \texttt{full} se nesmí prohodit.}

\ulohahead{společná informatika}{Objektové programování a návrh (C\#)}
\begin{zadani}
\begin{enumerate}[label=(\alph*)]
  \item \bod{1} Vysvětlete abstrakci, zapouzdření a polymorfismus a související konstrukty (třída, rozhraní, dědičnost, viditelnost). Rozdíl mezi statickým a dynamickým typováním a mezi virtuální a nevirtuální metodou.
  \item \bod{2} Navrhněte v C\# malou hierarchii geometrických tvaru s rozhraním a polymorfním výpočtem obsahu; ukažte ošetření chyby výjimkou.
  \item \bod{3} Jak je implementován dynamický polymorfismus (tabulka virtuálních metod)? K čemu jsou generické typy a jaký je rozdíl mezi hodnotovými a referenčními typy v C\#?
\end{enumerate}
\end{zadani}

\begin{reseni}
\textbf{(a)} \emph{Abstrakce} skrývá detaily za rozhraním; \emph{zapouzdření} drží data + operace pohromadě a omezuje přístup (viditelnost \texttt{private}/\texttt{public}/\texttt{protected}); \emph{polymorfismus} = jednotné volání nad různými typy. \emph{Třída} definuje typ, \emph{rozhraní} (interface) jen kontrakt, \emph{dědičnost} přebírá a rozšiřuje. \emph{Statické} typování kontroluje typy při překladu, \emph{dynamické} za běhu. \emph{Virtuální} metoda se vybírá podle skutečného typu objektu za běhu (override), \emph{nevirtuální} podle deklarovaného typu.

\textbf{(b)}
\begin{lstlisting}
public interface IShape { double Area(); }

public sealed class Circle : IShape {
    private readonly double r;
    public Circle(double r) {
        if (r < 0) throw new ArgumentException("zaporny polomer");
        this.r = r;
    }
    public double Area() => Math.PI * r * r;
}

public sealed class Rectangle : IShape {
    private readonly double w, h;
    public Rectangle(double w, double h) { this.w = w; this.h = h; }
    public double Area() => w * h;
}

double TotalArea(IEnumerable<IShape> shapes) {
    double sum = 0;
    foreach (var s in shapes) sum += s.Area(); // polymorfni volani
    return sum;
}
\end{lstlisting}

\textbf{(c)} \emph{Dynamický polymorfismus} se realizuje \emph{tabulkou virtuálních metod} (vtable): každý objekt drží ukazatel na tabulku své třídy, volání virtuální metody = skok přes položku tabulky (vybere se override potomka). \emph{Generické typy} umožní psát kód nezávislý na konkrétním typu (\texttt{List<T>}) bez ztráty typové kontroly a bez boxingu u hodnotových typu. \emph{Hodnotové} typy (\texttt{struct}, \texttt{int}) se kopírují podle hodnoty (typicky na zásobníku), \emph{referenční} (\texttt{class}) žijí na haldě a předává se reference; spravuje je garbage collector.
\end{reseni}

\kalibrace{Programovací jazyky / OO návrh je v moderním formátu vzácnější, ale když padne, bývá implementační. Definice (a) + funkční kostra hierarchie (b) tě udrží nad nulou; jazyk je C\# (volba majitele).}
\past{Rozhraní není abstraktní třída (nemá stav/implementaci, lze jich implementovat víc). Virtuální dispatch jde podle \emph{skutečného} typu - jinak by polymorfismus nefungoval.}


\section{Specializace: Strojové učení (Téma 3)}
\ulohahead{specializace UI-SU}{Lineární a logistická regrese, regularizace}
\begin{zadani}
\begin{enumerate}[label=(\alph*)]
  \item \bod{1} Definujte model lineární regrese a chybovou funkci (MSE). Napište \emph{normální rovnice} a analytické řešení metodou nejmenších čtverců.
  \item \bod{2} Co je L2 (ridge) regularizace? Jak změní normální rovnice a jaký má vliv na koeficienty? Proč pomáhá proti přeučení a kdy je nutná (singularita $X^{\!\top}X$)?
  \item \bod{3} Logistická regrese pro binární klasifikaci: napište predikci (sigmoid), ztrátovou funkci (NLL / křížová entropie) a gradient pro jeden krok SGD. Nemusíte nic počítat numericky.
\end{enumerate}
\end{zadani}

\begin{reseni}
\textbf{(a) Lineární regrese.} Data $X\in\mathbb{R}^{n\times d}$ (řádky $x_i$), cíle $y\in\mathbb{R}^n$, model $\hat y=Xw$ (sloupec jedniček v $X$ dává bias). Chyba:
\[
\mathrm{MSE}(w)=\tfrac1n\sum_{i=1}^n (x_i^{\!\top}w-y_i)^2=\tfrac1n\lVert Xw-y\rVert^2.
\]
Minimum: $\nabla_w=\tfrac2n X^{\!\top}(Xw-y)=0$, tj.\ \emph{normální rovnice}
\[
X^{\!\top}X\,w = X^{\!\top}y \quad\Rightarrow\quad w=(X^{\!\top}X)^{-1}X^{\!\top}y \ \text{(pokud je $X^{\!\top}X$ regulární).}
\]

\textbf{(b) L2 / ridge.} Minimalizujeme $\lVert Xw-y\rVert^2+\lambda\lVert w\rVert^2$, $\lambda>0$. Normální rovnice se změní na
\[
(X^{\!\top}X+\lambda I)\,w=X^{\!\top}y,\qquad w=(X^{\!\top}X+\lambda I)^{-1}X^{\!\top}y.
\]
Matice $X^{\!\top}X+\lambda I$ je vždy regulární (pozitivně definitní), takže řešení existuje i při kolineárních příznacích nebo $d>n$, kdy je $X^{\!\top}X$ singulární. Regularizace \emph{smršťuje} koeficienty směrem k nule (snižuje rozptyl modelu za cenu malého vychýlení), čímž omezuje přeučení. Větší $\lambda$ = silnější smrštění; $\lambda$ se volí křížovou validací.

\textbf{(c) Logistická regrese.} Predikce pravděpodobnosti třídy 1:
\[
p_i=\sigma(w^{\!\top}x_i)=\frac{1}{1+e^{-w^{\!\top}x_i}}.
\]
Ztráta (negativní log-věrohodnost / binární křížová entropie):
\[
L(w)=-\sum_{i=1}^n\big[y_i\ln p_i+(1-y_i)\ln(1-p_i)\big].
\]
Gradient má jednoduchý tvar $\nabla_w L=\sum_i (p_i-y_i)\,x_i$. Krok SGD na jednom příkladu $(x_i,y_i)$:
\[
w\leftarrow w-\eta\,(p_i-y_i)\,x_i,
\]
kde $\eta$ je rychlost učení. (Pro více tříd: softmax místo sigmoidu, ztráta = kategoriální křížová entropie, gradient $(p_i-\mathbb{1}_{y_i})x_i$.)
\end{reseni}

\begin{jadro}
Lineární: normální rovnice $X^{\!\top}X\,w=X^{\!\top}y$. Ridge: $(X^{\!\top}X+\lambda I)\,w=X^{\!\top}y$ (smršťuje, vždy řešitelné). Logistická: $p=\sigma(w^{\!\top}x)$, ztráta NLL, gradient $\sum_i(p_i-y_i)x_i$, SGD krok $w\leftarrow w-\eta(p_i-y_i)x_i$.
\end{jadro}

\kalibrace{Regrese je nejčastější ML okruh (~67\,\% zkoušek) a sytí kritickou podlahu specializace ($\ge 7/12$, kde potřebuješ tři ze čtyř otázek na $\ge 2$ body). Sigmoid, NLL a normální rovnice umět zpaměti = spolehlivé 2-3 body.}
\past{,,Metoda nejmenších čtverců'' = analytické řešení přes normální rovnice; SGD je iterativní alternativa. Pleteš-li si je, ztrácíš bod. U ridge nezapomeň, že bias ($w_0$) se obvykle neregularizuje.}

\ulohahead{specializace UI-SU}{Shlukování: k-means a hierarchické shlukování}
\begin{zadani}
\begin{enumerate}[label=(\alph*)]
  \item \bod{1} Definujte úlohu shlukování (učení bez učitele) a algoritmus k-means (krok přiřazení a krok aktualizace) a jeho účelovou funkci.
  \item \bod{2} Co je inicializace k-means++? Proč ztráta v k-means monotónně klesá a proč algoritmus konverguje (do lokálního minima)? Popište aglomerativní hierarchické shlukování a dendrogram.
  \item \bod{3} Jak zvolit počet shluku $k$? Proč je k-means citlivý na škálování příznaku? Proveďte jeden krok k-means pro body $1,2,\;10,11$ na přímce s centry $c_1=0,c_2=20$.
\end{enumerate}
\end{zadani}

\begin{reseni}
\textbf{(a)} \emph{Shlukování}: rozděl neoznačená data do $k$ skupin podle podobnosti. \emph{k-means} minimalizuje vnitroshlukový součet čtvercu $J=\sum_{i}\lVert x_i-\mu_{c(i)}\rVert^2$ střídáním: (1) \emph{přiřazení} každého bodu k nejbližšímu centru; (2) \emph{aktualizace} center na průměr přiřazených bodu.

\textbf{(b)} \emph{k-means++}: počáteční centra volí postupně s pravděpodobností úměrnou $D(x)^2$ (vzdálenost k nejbližšímu už zvolenému centru), což snižuje riziko špatného startu. Oba kroky $J$ nezvyšují (přiřazení k nejbližšímu i průměr jsou optimální dílčí kroky), $J$ je zdola omezené, konfigurací je konečně mnoho $\Rightarrow$ konvergence, ale jen do \emph{lokálního} minima (proto víc restartu). \emph{Aglomerativní} shlukování začíná s každým bodem zvlášť a opakovaně spojuje dva nejbližší shluky (linkage: single/complete/average); historii zachytí \emph{dendrogram}, který se ,,uřízne'' na požadovaný počet shluku.

\textbf{(c)} $k$ se volí např.\ metodou lokte (zlom v poklesu $J$) nebo silhouette skórem. \emph{Škálování}: euklidovská vzdálenost míchá jednotky, takže příznak s velkým rozsahem dominuje; proto se data standardizují. Krok pro body $\{1,2,10,11\}$, $c_1=0,c_2=20$: přiřazení: $1,2,10,11$ jsou blíž $c_1$ než $c_2$? $|10-0|=|10-20|=10$ (remíza), $|11-0|=11>|11-20|=9$, takže 11 jde k $c_2$. Tedy $\{1,2,10\}\to c_1$, $\{11\}\to c_2$ (remízu u 10 dáme $c_1$). Aktualizace: $c_1=\tfrac{1+2+10}{3}=4{,}33$, $c_2=11$. (Další iterace už rozdělí na $\{1,2\}$ a $\{10,11\}$.)
\end{reseni}

\kalibrace{Shlukování je druhý nejčastější ML okruh (~44\,\%, k-means se opakuje). Definice + jeden krok výpočtu = 2 body do kritické podlahy specializace.}
\past{k-means najde jen \emph{lokální} minimum a předpokládá zhruba kulové shluky; není deterministický (závisí na inicializaci). Před během škáluj příznaky.}

\ulohahead{specializace UI-SU}{Evaluační metriky binárního klasifikátoru}
\begin{zadani}
\begin{enumerate}[label=(\alph*)]
  \item \bod{1} Nakreslete matici záměn a definujte accuracy, precision, recall a F1.
  \item \bod{2} Proč je accuracy zavádějící u nevyvážených tříd? Kdy upřednostnit precision a kdy recall? Z hodnot $\mathrm{TPR}$, $\mathrm{FPR}$ a počtu pozitivních/negativních rekonstruujte $\mathrm{TP,FP,TN,FN}$.
  \item \bod{3} Co zobrazuje ROC křivka a co je AUC? Spočtěte precision a recall pro $\mathrm{TP}=20,\ \mathrm{FP}=30,\ \mathrm{FN}=5$.
\end{enumerate}
\end{zadani}

\begin{reseni}
\textbf{(a)} Matice záměn (řádky skutečnost, sloupce predikce):
\[
\begin{array}{c|cc}
 & \hat{+} & \hat{-}\\\hline
+ & \mathrm{TP} & \mathrm{FN}\\
- & \mathrm{FP} & \mathrm{TN}
\end{array}
\]
$\text{accuracy}=\frac{\mathrm{TP+TN}}{\text{vše}}$, $\text{precision}=\frac{\mathrm{TP}}{\mathrm{TP+FP}}$, $\text{recall (TPR)}=\frac{\mathrm{TP}}{\mathrm{TP+FN}}$, $F_1=\frac{2\,PR}{P+R}$ (harmonický průměr).

\textbf{(b)} Při nevyvážených třídách (např.\ 1\,\% pozitivních) dá triviální klasifikátor ,,vše negativní'' accuracy 99\,\%, ač nic nenajde - proto se sleduje precision/recall. \emph{Recall} je důležitý, když je drahé minout pozitivní (nemoc, podvod); \emph{precision}, když je drahý falešný poplach (spam označený jako důležitý). Rekonstrukce z $\mathrm{TPR},\mathrm{FPR}$ a počtu $P$ (pozitivních) a $N$ (negativních): $\mathrm{TP}=\mathrm{TPR}\cdot P$, $\mathrm{FN}=P-\mathrm{TP}$, $\mathrm{FP}=\mathrm{FPR}\cdot N$, $\mathrm{TN}=N-\mathrm{FP}$.

\textbf{(c)} \emph{ROC} vynáší $\mathrm{TPR}$ proti $\mathrm{FPR}$ při měnícím se prahu klasifikace; \emph{AUC} (plocha pod ní) je pravděpodobnost, že náhodný pozitivní dostane vyšší skóre než náhodný negativní (1 = ideál, 0{,}5 = náhoda). Pro zadání: $\text{precision}=\frac{20}{20+30}=0{,}4$, $\text{recall}=\frac{20}{20+5}=0{,}8$.
\end{reseni}

\begin{jadro}
$\text{precision}=\frac{TP}{TP+FP}$, $\text{recall}=\frac{TP}{TP+FN}$, $F_1=\frac{2PR}{P+R}$ (harmonický). U nevyvážených tříd accuracy klame. Rekonstrukce: $TP=\mathrm{TPR}\cdot P$, $FP=\mathrm{FPR}\cdot N$. ROC = TPR vs FPR.
\end{jadro}

\kalibrace{Evaluační metriky jsou velmi častý ML okruh (~44\,\%) a opakovaně chtějí rekonstrukci počtu z metrik. Vzorce P/R/F1 a převod TPR/FPR umět zpaměti = spolehlivé body.}
\past{Precision a recall si nepleť (jmenovatel: u precision predikované pozitivní, u recall skutečně pozitivní). F1 je harmonický, ne aritmetický průměr.}

\ulohahead{specializace UI-SU}{Rozhodovací stromy: kritéria větvení a prořezávání}
\begin{zadani}
\begin{enumerate}[label=(\alph*)]
  \item \bod{1} Definujte rozhodovací strom a kritéria větvení: entropii a informační zisk, Gini index.
  \item \bod{2} Popište hladovou konstrukci stromu. Spočtěte informační zisk štěpení, které z uzlu se 4 kladnými a 4 zápornými příklady udělá dva čisté listy (4/0 a 0/4).
  \item \bod{3} Co je přeučení u stromu a jak mu brání prořezávání (pre-pruning vs post-pruning)?
\end{enumerate}
\end{zadani}

\begin{reseni}
\textbf{(a)} \emph{Rozhodovací strom}: vnitřní uzly testují příznak, listy přiřazují třídu (či hodnotu). \emph{Entropie} množiny $S$ s podíly tříd $p_k$: $H(S)=-\sum_k p_k\log_2 p_k$. \emph{Informační zisk} štěpení $S$ na části $S_j$: $\mathrm{IG}=H(S)-\sum_j\frac{|S_j|}{|S|}H(S_j)$. \emph{Gini}: $G(S)=1-\sum_k p_k^2$ (alternativní míra nečistoty).

\textbf{(b)} \emph{Konstrukce} (ID3/CART): v každém uzlu zvol příznak s největším informačním ziskem (resp.\ nejmenší váženou nečistotou), rozděl data a rekurzivně pokračuj, dokud nejsou listy čisté nebo nenastane zastavení. Výpočet: rodič $4{+}/4{-}$ má $H=-\tfrac12\log_2\tfrac12-\tfrac12\log_2\tfrac12=1$ bit. Děti $4/0$ a $0/4$ jsou čisté: $H=0$. $\mathrm{IG}=1-\big(\tfrac12\cdot0+\tfrac12\cdot0\big)=1$ bit (maximální možný zisk - dokonalé rozdělení).

\textbf{(c)} \emph{Přeučení}: strom roste do hloubky, učí se šum a na trénovacích datech má nulovou chybu, ale špatně generalizuje. \emph{Pre-pruning} (předčasné zastavení): omezení hloubky, minimální počet vzorku v listu, minimální zisk. \emph{Post-pruning}: nech strom dorůst a pak ořezávej podstromy, které nezlepšují chybu na validační množině (nahrazení listem).
\end{reseni}

\kalibrace{Rozhodovací stromy jsou častý okruh (~33\,\%) a vděčí za body za entropii/Gini a jeden výpočet IG. Patří do plně zvládnutého jádra specializace.}
\past{Entropie se počítá s $\log_2$ (bity). Informační zisk je \emph{vážený} podle velikostí dětí, ne prostý rozdíl entropií.}

\ulohahead{specializace UI-SU}{Kombinace modelu: bagging, boosting, náhodné lesy}
\begin{zadani}
\begin{enumerate}[label=(\alph*)]
  \item \bod{1} Co je kombinace více modelu (ensemble)? Vysvětlete rozdíl mezi baggingem a boostingem.
  \item \bod{2} Popište náhodný les (random forest) a vysvětlete, proč zlepšuje přesnost oproti jednomu stromu.
  \item \bod{3} Popište princip AdaBoostu. Za jakých podmínek ensemble pomáhá nejvíc a kdy vůbec ne (majoritní hlasování)?
\end{enumerate}
\end{zadani}

\begin{reseni}
\textbf{(a)} \emph{Ensemble} kombinuje predikce více modelu (hlasování/průměr) do silnějšího. \emph{Bagging} trénuje modely \emph{nezávisle} na bootstrap vzorcích a průměruje (snižuje hlavně \emph{rozptyl}). \emph{Boosting} trénuje modely \emph{postupně}, každý se zaměří na chyby předchozích (snižuje hlavně \emph{vychýlení}).

\textbf{(b)} \emph{Random forest} = bagging rozhodovacích stromu, kde se navíc v každém štěpení uvažuje jen náhodná podmnožina příznaku. To stromy \emph{dekoreluje}; protože průměr nekorelovaných chyb má menší rozptyl, les generalizuje lépe než jeden hluboký (přeučený) strom, a přitom téměř nepotřebuje ladění.

\textbf{(c)} \emph{AdaBoost}: udržuje váhy příkladu, v každém kole natrénuje slabý klasifikátor, zvýší váhy špatně klasifikovaných a přidá klasifikátor s váhou podle jeho přesnosti; finální predikce = vážené hlasování. Ensemble pomáhá nejvíc, jsou-li dílčí modely \emph{lepší než náhoda} a dělají \emph{různé/nekorelované} chyby. Naopak nepomůže, dělají-li všechny stejné chyby (dokonale korelované) - pak majoritní hlasování dopadne stejně jako jeden model; nejlepší případ je, když se chyby nepřekrývají a většina vždy ,,přehlasuje'' menšinu chybujících.
\end{reseni}

\begin{jadro}
Bagging = trénuj nezávisle na bootstrap vzorcích + průměruj (snižuje rozptyl). Boosting = sekvenčně, každý model na chybách předchozích (snižuje vychýlení). Random forest = bagging stromu + náhodná podmnožina příznaku (dekorelace). Pomáhá při nekorelovaných chybách.
\end{jadro}

\kalibrace{Kombinace modelu (~22\,\%) je vděčná konceptuální otázka (žádný výpočet). Bagging vs boosting + random forest = 2-3 body za vysvětlení.}
\past{Bagging $\neq$ boosting: bagging paralelně/nezávisle (rozptyl), boosting sekvenčně na chybách (vychýlení). Les pomáhá díky \emph{dekorelaci} stromu, ne jen jejich počtu.}

\ulohahead{specializace UI-SU}{Metoda podpůrných vektoru (SVM)}
\begin{zadani}
\begin{enumerate}[label=(\alph*)]
  \item \bod{1} Popište SVM pro lineárně separabilní třídy: co je rozdělující nadrovina, margin a podpůrné vektory (support vectors)?
  \item \bod{2} Co je měkký margin a parametr $C$? Jaký je kompromis mezi šířkou marginu a počtem chyb?
  \item \bod{3} K čemu slouží jádrové funkce (kernel trick) a jak se chová RBF jádro? Jak SVM klasifikuje do více tříd?
\end{enumerate}
\end{zadani}

\begin{reseni}
\textbf{(a)} SVM hledá \emph{rozdělující nadrovinu} $w^{\!\top}x+b=0$ s \emph{maximálním marginem} (pásem) mezi třídami. Margin je vzdálenost nadroviny k nejbližším bodum; ty body leží na okraji a nazývají se \emph{podpůrné vektory} - jen ony určují řešení (ostatní lze odebrat beze změny). Maximalizace marginu $=$ minimalizace $\tfrac12\lVert w\rVert^2$ za podmínek $y_i(w^{\!\top}x_i+b)\ge1$.

\textbf{(b)} Když data nejsou separabilní, zavede se \emph{měkký margin} s prom.\ uvolnění $\xi_i\ge0$: minimalizujeme $\tfrac12\lVert w\rVert^2+C\sum_i\xi_i$. Parametr $C$ váží chyby proti šířce marginu: \emph{velké} $C$ = málo chyb, úzký margin (riziko přeučení); \emph{malé} $C$ = širší margin, toleruje víc chyb (lepší generalizace).

\textbf{(c)} \emph{Jádrový trik}: nahradí skalární součin $x_i^{\!\top}x_j$ jádrem $K(x_i,x_j)$, čímž SVM počítá v (implicitně vysokorozměrném) prostoru rysu, aniž ho explicitně sestrojí - umožní nelineární hranice. \emph{RBF} $K(x,z)=\exp(-\gamma\lVert x-z\rVert^2)$ měří podobnost lokálně; malé $\gamma$ = hladká hranice, velké $\gamma$ = klikatá (přeučení). \emph{Více tříd}: kombinací binárních klasifikátoru one-vs-rest nebo one-vs-one.
\end{reseni}

\kalibrace{SVM (~12\,\%) je vzácnější, ale když padne, je konceptuální (,,popište, nakreslete hranici, vliv bodu, kdy RBF''). Definice marginu + role $C$/jádra = 2-bodová záloha.}
\past{Hranici určují jen podpůrné vektory. Velké $C$ = méně tolerance chyb (užší margin), ne naopak; RBF s velkým $\gamma$ se přeučuje.}

\ulohahead{specializace UI-SU}{Analýza hlavních komponent (PCA)}
\begin{zadani}
\begin{enumerate}[label=(\alph*)]
  \item \bod{1} Co je PCA a k čemu slouží? Co jsou hlavní komponenty a jak souvisí s kovarianční maticí?
  \item \bod{2} Popište postup PCA (centrování, kovariance, vlastní rozklad / SVD, projekce na top-$k$). Co je vysvětlený rozptyl?
  \item \bod{3} Jaká je souvislost PCA a SVD? Proč je nutné data před PCA škálovat a co PCA optimalizuje (rekonstrukční chyba)?
\end{enumerate}
\end{zadani}

\begin{reseni}
\textbf{(a)} \emph{PCA} je lineární redukce dimenze: hledá nové ortogonální osy (\emph{hlavní komponenty}), podél nichž mají data největší rozptyl, a projektuje na prvních $k$ z nich. Hlavní komponenty jsou \emph{vlastní vektory kovarianční matice} dat; příslušná vlastní čísla udávají rozptyl podél nich.

\textbf{(b)} Postup: (1) \emph{centruj} data (odečti průměr), případně standardizuj; (2) spočti \emph{kovarianční matici} $C=\tfrac1n X^{\!\top}X$; (3) najdi její \emph{vlastní čísla a vektory} (nebo SVD matice $X$); (4) seřaď komponenty podle vlastních čísel a \emph{projektuj} na prvních $k$. \emph{Vysvětlený rozptyl} komponenty = $\lambda_i/\sum_j\lambda_j$; volíme $k$ tak, aby pokryl např.\ 95\,\% rozptylu.

\textbf{(c)} Pro centrovaná data $X=U\Sigma V^{\!\top}$ (SVD) jsou hlavní komponenty sloupce $V$ a singulární hodnoty $\sigma_i$ souvisí s vlastními čísly kovariance přes $\lambda_i=\sigma_i^2/n$; SVD je numericky stabilnější cesta. \emph{Škálování} je nutné, protože PCA reaguje na rozptyl: příznak s velkou jednotkou (rozsahem) by jinak uměle dominoval první komponentě. PCA minimalizuje \emph{rekonstrukční chybu} (součet čtvercu kolmých vzdáleností k podprostoru), což je ekvivalentní maximalizaci zachyceného rozptylu.
\end{reseni}

\kalibrace{PCA (~12\,\%) je vzácnější, ale opakované jako konceptuální + vizuální otázka. Definice (komponenty = vlastní vektory kovariance) + postup = 2-bodová záloha.}
\past{PCA hledá směry maximálního \emph{rozptylu}, ne směry oddělující třídy (to je LDA). Bez centrování/škálování vyjdou komponenty zkreslené.}

\ulohahead{specializace UI-SU}{k-nejbližších sousedu (kNN)}
\begin{zadani}
\begin{enumerate}[label=(\alph*)]
  \item \bod{1} Popište algoritmus k-nejbližších sousedu pro klasifikaci i regresi. Proč se mu říká ,,líné'' učení (lazy)?
  \item \bod{2} Jak se volí $k$ (křížová validace) a jak malé vs velké $k$ ovlivní výsledek (vychýlení vs rozptyl)? Proč je nutné škálovat příznaky?
  \item \bod{3} Co je prokletí dimenzionality a jak postihuje kNN? Jaká je výpočetní složitost predikce?
\end{enumerate}
\end{zadani}

\begin{reseni}
\textbf{(a)} kNN si \emph{pamatuje} trénovací data; pro nový bod najde $k$ nejbližších (dle metriky, typicky euklidovské) a \emph{klasifikace} = většinové hlasování jejich tříd, \emph{regrese} = průměr jejich hodnot. Je \emph{líné}, protože netrénuje žádný model - veškerá práce je až při predikci.

\textbf{(b)} $k$ se volí \emph{křížovou validací}. \emph{Malé} $k$ (např.\ 1) dává nízké vychýlení, ale vysoký rozptyl (citlivé na šum, přeučení); \emph{velké} $k$ hranici vyhladí (vyšší vychýlení, nižší rozptyl). \emph{Škálování} je nutné, protože vzdálenost jinak ovládne příznak s největším rozsahem.

\textbf{(c)} \emph{Prokletí dimenzionality}: ve vysoké dimenzi se vzdálenosti mezi body vyrovnávají (poměr nejbližší/nejvzdálenější $\to1$), takže pojem ,,soused'' ztrácí smysl a kNN selhává; navíc je třeba exponenciálně víc dat. \emph{Složitost} naivní predikce je $O(nd)$ na jeden dotaz (projít $n$ bodu v dimenzi $d$); urychlení přes k-d stromy či přibližné hledání.
\end{reseni}

\begin{jadro}
kNN = ulož data, pro nový bod najdi $k$ nejbližších; klasifikace = hlasování, regrese = průměr. Líné učení. $k$ přes křížovou validaci (malé $k$ = rozptyl, velké = vychýlení). Škáluj příznaky; trpí prokletím dimenzionality.
\end{jadro}

\kalibrace{kNN (~12\,\%) je vzácnější, ale jednoduchý a vděčný. Definice + role $k$ a škálování = spolehlivá 2-bodová záloha; navíc nese pojem prokletí dimenzionality.}
\past{kNN netrénuje (lazy) - ,,trénink'' je jen uložení dat. Bez škálování dominuje příznak s velkou jednotkou; pozor i na liché $k$ u dvou tříd (remízy).}

\ulohahead{specializace UI-SU}{Statistické testy: t-test a chí-kvadrát}
\begin{zadani}
\begin{enumerate}[label=(\alph*)]
  \item \bod{1} Co je statistický test, nulová hypotéza, hladina významnosti a p-hodnota? Popište jednovýběrový a dvouvýběrový t-test.
  \item \bod{2} Popište chí-kvadrát test dobré shody: testovou statistiku, stupně volnosti a rozhodovací pravidlo.
  \item \bod{3} Co je párový t-test a kdy ho použít? Jak souvisí rozhodnutí testu s chybami 1.\ a 2.\ druhu?
\end{enumerate}
\end{zadani}

\begin{reseni}
\textbf{(a)} \emph{Test} rozhoduje mezi \emph{nulovou hypotézou} $H_0$ a alternativou na základě dat. \emph{Hladina významnosti} $\alpha$ = přípustná pravděpodobnost chyby 1.\ druhu; \emph{p-hodnota} = pravděpodobnost dat aspoň tak extrémních jako pozorovaná, platí-li $H_0$; zamítáme, je-li $p<\alpha$. \emph{Jednovýběrový} t-test: $H_0:\mu=\mu_0$, statistika $t=\frac{\bar x-\mu_0}{s/\sqrt n}$, která má za platnosti $H_0$ Studentovo $t$-rozdělení s $n-1$ stupni volnosti; zamítáme, padne-li $t$ do kritického oboru (resp.\ $p<\alpha$). \emph{Dvouvýběrový}: $H_0:\mu_1=\mu_2$, statistika $t=\frac{\bar x-\bar y}{\sqrt{s_X^2/n+s_Y^2/m}}$ (Welchova varianta při různých rozptylech).

\textbf{(b)} \emph{Chí-kvadrát test dobré shody} porovná pozorované četnosti $O_i$ s očekávanými $E_i$ podle modelu: $\chi^2=\sum_i\frac{(O_i-E_i)^2}{E_i}$. \emph{Stupně volnosti} $=$ (počet kategorií $-1-$ počet odhadovaných parametru). $H_0$ (data odpovídají modelu) zamítneme, je-li $\chi^2$ větší než kritická hodnota $\chi^2_{1-\alpha}$ pro dané df.

\textbf{(c)} \emph{Párový} t-test se použije, jsou-li data \emph{spárovaná} (např.\ měření před/po na týchž jedincích); testuje, zda je průměr rozdílu nulový. \emph{Chyba 1.\ druhu} = zamítnout platné $H_0$ (řízena $\alpha$); \emph{chyba 2.\ druhu} = nezamítnout neplatné $H_0$ (řízena silou testu $1-\beta$, roste s velikostí vzorku a efektu).
\end{reseni}

\begin{jadro}
Test: $H_0$, testová statistika, p-hodnota; zamítni při $p<\alpha$. Jednovýběrový t-test $t=\frac{\bar x-\mu_0}{s/\sqrt n}$. Chí-kvadrát dobré shody $\chi^2=\sum_i\frac{(O_i-E_i)^2}{E_i}$, df $=$ kategorie $-1-$ parametry.
\end{jadro}

\kalibrace{Statistické testy (~22\,\% dohromady) chtějí převážně definice + dosazení do statistiky. Stačí 2-bodová záloha: hypotéza, statistika, rozhodnutí.}
\past{p-hodnota není pravděpodobnost, že $H_0$ platí. Stupně volnosti u chí-kvadrát nejsou počet kategorií, ale po odečtení vazeb/parametru.}

\ulohahead{specializace UI-SU}{Vyhodnocení modelu: validace, přeučení, regularizace}
\begin{zadani}
\begin{enumerate}[label=(\alph*)]
  \item \bod{1} Vysvětlete rozdělení dat na train/dev/test a křížovou validaci ($k$-fold). Co je princip maximální věrohodnosti?
  \item \bod{2} Co je přeučení (overfitting) a generalizační chyba? Jak proti přeučení působí regularizace (L1 vs L2) a včasné zastavení trénování?
  \item \bod{3} Vysvětlete kompromis vychýlení-rozptyl (bias-variance) a prokletí dimenzionality při výběru modelu.
\end{enumerate}
\end{zadani}

\begin{reseni}
\textbf{(a)} \emph{Train} slouží k učení, \emph{dev/validace} k ladění hyperparametru a výběru modelu, \emph{test} k jedinému závěrečnému odhadu výkonu (nesmí se na něm ladit). \emph{$k$-fold křížová validace}: data rozděl na $k$ částí, $k$-krát trénuj na $k{-}1$ a validuj na zbylé, výsledky zprůměruj (lepší využití dat). \emph{Maximální věrohodnost}: zvol parametry maximalizující pravděpodobnost (věrohodnost) pozorovaných dat, $\hat\theta=\arg\max_\theta P(\text{data}\mid\theta)$.

\textbf{(b)} \emph{Přeučení}: model se naučí šum trénovacích dat, má malou trénovací, ale velkou \emph{generalizační} (testovací) chybu. \emph{Regularizace} přidá penaltu za velikost vah: \emph{L2} ($\lambda\lVert w\rVert_2^2$) hladce smršťuje, \emph{L1} ($\lambda\lVert w\rVert_1$) tlačí váhy přesně na nulu (řídké řešení, výběr příznaku). \emph{Včasné zastavení}: trénink se ukončí, když začne růst validační chyba.

\textbf{(c)} \emph{Bias-variance}: příliš jednoduchý model má velké \emph{vychýlení} (podučení), příliš složitý velký \emph{rozptyl} (přeučení); cílem je minimum jejich součtu. \emph{Prokletí dimenzionality}: s rostoucím počtem příznaku roste potřeba dat exponenciálně a vzdálenosti se vyrovnávají, takže složitější modely snáz přeučí - proto redukce dimenze/regularizace.
\end{reseni}

\begin{jadro}
Train (učení) / dev (ladění) / test (jediný závěrečný odhad); $k$-fold CV průměruje. Přeučení = malá train, velká test chyba. L2 hladce smršťuje, L1 dělá řídké váhy (výběr příznaku). Bias-variance: minimalizuj součet.
\end{jadro}

\kalibrace{Toto je ,,zastřešující'' ML téma (evaluace, regularizace) prolínající mnoho otázek; opakovaně se ptá na křížovou validaci a přeučení. Solidní definice = body napříč specializací.}
\past{Na testovacích datech se \emph{neladí}. L1 dělá řídké váhy (výběr příznaku), L2 jen smršťuje - nezaměňuj jejich efekt.}

\ulohahead{specializace UI-SU}{Predikce z tabulárních dat: předzpracování a metriky}
\begin{zadani}
\begin{enumerate}[label=(\alph*)]
  \item \bod{1} Popište předzpracování tabulárních dat: kódování kategoriálních příznaku (one-hot), škálování spojitých příznaku a ošetření chybějících hodnot (imputace).
  \item \bod{2} Definujte regresní metriky MAE, RMSE a $R^2$. Kdy je vhodnější MAE a kdy RMSE?
  \item \bod{3} Jaký model zvolit pro tabulární regresi (lineární / stromové / náhodný les / boosting) a proč stromové metody nepotřebují škálování?
\end{enumerate}
\end{zadani}

\begin{reseni}
\textbf{(a)} \emph{Kategoriální} příznaky se převedou na čísla: \emph{one-hot} kódování (každá hodnota = binární sloupec) pro neuspořádané kategorie; ordinální kódování pro uspořádané. \emph{Spojité} příznaky se \emph{škálují} (standardizace na nulový průměr a jednotkový rozptyl, nebo min-max) - nutné pro vzdálenostní a gradientové modely. \emph{Chybějící hodnoty}: imputace (průměr/medián/nejčastější hodnota), případně příznak ,,bylo chybějící''.

\textbf{(b)} Pro predikce $\hat y_i$ a cíle $y_i$: $\mathrm{MAE}=\tfrac1n\sum|y_i-\hat y_i|$; $\mathrm{RMSE}=\sqrt{\tfrac1n\sum(y_i-\hat y_i)^2}$; $R^2=1-\frac{\sum(y_i-\hat y_i)^2}{\sum(y_i-\bar y)^2}$ (podíl vysvětleného rozptylu). \emph{MAE} je odolnější vuči odlehlým hodnotám; \emph{RMSE} velké chyby trestá kvadraticky (preferuj, když jsou velké chyby zvlášť nežádoucí).

\textbf{(c)} \emph{Lineární regrese} pro zhruba lineární vztahy (rychlá, interpretovatelná); \emph{stromové metody} (rozhodovací strom, \emph{náhodný les}, \emph{gradient boosting}) pro nelineární vztahy a smíšené typy příznaku - na tabulárních datech bývají nejlepší a robustní. Stromy dělí podle \emph{prahu} jednoho příznaku, takže jsou \emph{invariantní k monotónní transformaci} (škálování nezmění pořadí hodnot, tedy ani štěpení) - proto nepotřebují normalizaci.
\end{reseni}

\kalibrace{Tabulární workflow (~10\,\%) se objevil nedávno a chtěl předzpracování + metriky (ne jen definici modelu). MAE/RMSE a one-hot/škálování/imputace = 2-bodová záloha proti ,,end-to-end'' otázce.}
\past{One-hot pro neuspořádané kategorie (jinak modelu podsouváš falešné pořadí). RMSE je v jednotkách cíle (odmocnina), MSE ne; stromy škálování nepotřebují, lineární/kNN/SVM ano.}


\section{Specializace: Základy umělé inteligence (Téma 1)}
\ulohahead{specializace UI-SU}{Splňování podmínek (CSP), hranová konzistence, AC-3}
\begin{zadani}
\begin{enumerate}[label=(\alph*)]
  \item \bod{1} Definujte problém splňování podmínek (CSP) a pojem \emph{hranové konzistence} (arc consistency).
  \item \bod{2} Popište algoritmus AC-3 (fronta hran, revize domény, šíření) a uveďte jeho časovou složitost.
  \item \bod{3} Pro proměnné $X,Y,Z$ s doménami $\{1,2,3\}$ a omezeními $X<Y$, $Y<Z$ proveďte AC-3 a uveďte výsledné domény. Vysvětlete rozdíl mezi lokální (hranovou) a globální konzistencí a proč hranová konzistence nezaručí existenci řešení.
\end{enumerate}
\end{zadani}

\begin{reseni}
\textbf{(a) Definice.} CSP je trojice $(\mathcal{V},\mathcal{D},\mathcal{C})$: proměnné $\mathcal{V}=\{X_1,\dots,X_n\}$, jejich domény $\mathcal{D}=\{D_1,\dots,D_n\}$ a omezení $\mathcal{C}$ (každé omezuje přípustné kombinace hodnot podmnožiny proměnných). Řešení = přiřazení hodnot z domén splňující všechna omezení. Hrana $(X_i,X_j)$ (binární omezení) je \emph{hranově konzistentní}, pokud pro každou hodnotu $a\in D_i$ existuje hodnota $b\in D_j$ taková, že dvojice $(a,b)$ splňuje omezení. CSP je hranově konzistentní, jsou-li konzistentní všechny hrany.

\textbf{(b) AC-3.} Udržuj frontu všech hran (orientovaně, obě strany každého omezení). Opakovaně vyjmi hranu $(X_i,X_j)$ a proveď \textsc{Revise}: odstraň z $D_i$ každou hodnotu, pro kterou v $D_j$ neexistuje podpora. Pokud se $D_i$ zmenšila, vlož zpět do fronty všechny hrany $(X_k,X_i)$ pro sousedy $X_k$ (mohla se porušit jejich konzistence). Konec, když je fronta prázdná, nebo nějaká doména zůstane prázdná (pak CSP nemá řešení). Složitost: $O(e\,d^3)$, kde $e$ je počet hran (binárních omezení) a $d$ velikost domény (každá z $O(e)$ hran se zařadí nejvýše $d$-krát, jedna \textsc{Revise} je $O(d^2)$).

\textbf{(c) Průběh.} Z $X<Y$: $X$ nemůže být $3$ (nic většího v $Y$) $\Rightarrow D_X=\{1,2\}$; $Y$ nemůže být $1$ $\Rightarrow D_Y=\{2,3\}$. Z $Y<Z$: $Y$ nemůže být $3$ $\Rightarrow D_Y=\{2\}$; $Z$ nemůže být $1$ ani $2$ $\Rightarrow D_Z=\{3\}$. Zmenšení $D_Y$ vrátí do fronty hranu $(X,Y)$: $X$ nyní potřebuje $y>x$ s $y\in\{2\}$, takže $x=2$ vypadne $\Rightarrow D_X=\{1\}$. Výsledek: $D_X=\{1\},\,D_Y=\{2\},\,D_Z=\{3\}$ (zde jediné řešení).

\emph{Lokální vs.\ globální.} Hranová konzistence je \emph{lokální} vlastnost (po dvojicích). Nezaručuje globální řešení: např.\ tři proměnné s doménami $\{1,2\}$ a omezeními ,,navzájem různé'' jsou hranově konzistentní (každá hodnota má podporu u souseda), ale tři po dvou různé hodnoty z dvouprvkové domény neexistují, takže CSP je nesplnitelný. Proto se AC-3 kombinuje s prohledáváním (backtracking), typicky algoritmus MAC = po každém přiřazení znovu udělat hranovou konzistenci.
\end{reseni}

\begin{jadro}
CSP $=(V,D,C)$. Hranová konzistence: každá hodnota $a\in D_i$ má podporu $b\in D_j$ splňující omezení. AC-3 (fronta hran, revize, šíření při zúžení), $O(e\,d^3)$. Nezaručuje řešení -> nutný backtracking/MAC.
\end{jadro}

\kalibrace{Specializaci tvoří ~3 ML + 1 ,,Základy UI''. AI okruhy (CSP, minimax, Bayesovské sítě, A*, MDP, HMM) jsou jednotlivě vzácné, ale slot na jednu z nich je skoro vždy. Proto v úrovni 1 musí mít každá AI rodina aspoň 2-bodovou zálohu: přesná definice + jeden kanonický postup.}
\past{Hranová konzistence \emph{nedokazuje} řešitelnost - to je oblíbená chytací podotázka. Vždy zmiň protipříklad (např.\ all-different na malé doméně) a roli backtrackingu/MAC.}

\ulohahead{specializace UI-SU}{Prohledávání: neinformované a informované (A*)}
\begin{zadani}
\begin{enumerate}[label=(\alph*)]
  \item \bod{1} Popište formulaci úlohy prohledáváním (stav, operátory, cíl, cena). Rozdíl mezi stromovým a grafovým prohledáváním. Uveďte neinformované metody BFS, DFS a uniform-cost.
  \item \bod{2} Popište algoritmus A* a funkci $f(n)=g(n)+h(n)$. Co je přípustná a co konzistentní heuristika?
  \item \bod{3} Za jakých podmínek je A* optimální? Co znamená dominance heuristik a proč je lepší ,,větší'' přípustná heuristika?
\end{enumerate}
\end{zadani}

\begin{reseni}
\textbf{(a)} \emph{Formulace}: počáteční stav, množina \emph{operátoru} (akce s cenou) generujících následníky, \emph{cílový test}, cena cesty. \emph{Stromové} prohledávání neeviduje navštívené stavy (může opakovat), \emph{grafové} drží uzavřenou množinu (closed set). \emph{BFS}: po vrstvách (fronta), úplné, optimální při jednotkových cenách. \emph{DFS}: do hloubky (zásobník), paměťově levné, neúplné/nesuboptimální. \emph{Uniform-cost}: rozšiřuje uzel s nejmenší $g$ (Dijkstra), optimální pro nezáporné ceny.

\textbf{(b)} \emph{A*} rozšiřuje uzel s nejmenším $f(n)=g(n)+h(n)$, kde $g$ je cena z počátku a $h$ \emph{heuristický} odhad ceny do cíle. Heuristika je \emph{přípustná}, pokud nikdy nepřeceňuje ($0\le h(n)\le h^*(n)$); \emph{konzistentní} (monotónní), pokud $h(n)\le c(n,n')+h(n')$ pro každý přechod (trojúhelníková nerovnost).

\textbf{(c)} A* je optimální se \emph{přípustnou} heuristikou ve stromovém prohledávání a s \emph{konzistentní} heuristikou v grafovém (jinak je nutné znovuotevírání uzlu). \emph{Dominance}: $h_2$ dominuje $h_1$, je-li $h_2(n)\ge h_1(n)$ pro všechny $n$ (a obě přípustné); pak A* s $h_2$ rozšíří méně uzlu (efektivnější), protože vyšší přípustný odhad ořeže víc.
\end{reseni}

\kalibrace{Prohledávání (~15\,\%) je opakovaná ,,Základy UI'' otázka (A* se objevil). Definice A* + přípustnost/konzistence = 2-bodová záloha pro AI slot specializace.}
\past{Přípustná $\neq$ konzistentní (konzistence je silnější a implikuje přípustnost). V grafovém A* je pro optimalitu potřeba konzistence, ne jen přípustnost.}

\ulohahead{specializace UI-SU}{Hry: minimax, alfa-beta a teorie her}
\begin{zadani}
\begin{enumerate}[label=(\alph*)]
  \item \bod{1} Popište algoritmus Minimax pro hru dvou hráču s nulovým součtem. Co předpokládá o protihráči?
  \item \bod{2} Popište alfa-beta prořezávání: co a proč prořezává a jak pořadí tahu ovlivní úsporu.
  \item \bod{3} Vysvětlete základy teorie her: vězňovo dilema, dominantní strategie, Nashovo ekvilibrium a Paretovskou optimalitu. Co je mechanism design a jaké jsou typy aukcí?
\end{enumerate}
\end{zadani}

\begin{reseni}
\textbf{(a)} \emph{Minimax}: prohledá herní strom; hráč MAX volí tah maximalizující hodnotu, MIN minimalizující. Hodnota listu = užitek; hodnota vnitřního uzlu = max/min hodnot synu podle toho, kdo je na tahu. Předpokládá \emph{optimálně hrajícího} protihráče (nejhorší případ pro nás). Hloubku lze omezit a listy nahradit \emph{ohodnocovací funkcí}.

\textbf{(b)} \emph{Alfa-beta} udržuje $\alpha$ (nejlepší jistá hodnota pro MAX) a $\beta$ (pro MIN). Větev se \emph{ořízne}, jakmile $\alpha\ge\beta$, protože ji racionální hráč stejně nezvolí, takže nemá smysl ji dál prohledávat. Výsledek je identický s minimaxem. Při \emph{dobrém pořadí} tahu (nejlepší první) klesne počet prohledaných uzlu z $O(b^d)$ až na $O(b^{d/2})$, tj.\ dvojnásobná dosažitelná hloubka.

\textbf{(c)} \emph{Vězňovo dilema}: dva hráči, každý volí spolupráci/zradu; zrada je \emph{dominantní strategie} (lepší bez ohledu na soupeře), takže oba zradí, ač by oboustranná spolupráce byla lepší. \emph{Nashovo ekvilibrium}: profil strategií, kde nikdo nezíská jednostrannou změnou (vzájemně nejlepší odpovědi); ve vězňově dilematu je to (zrada, zrada). \emph{Paretovsky optimální} je výsledek, který nelze zlepšit jednomu bez zhoršení druhému - (spolupráce, spolupráce) je Pareto-lepší, ale není Nashovo ekvilibrium. \emph{Mechanism design} je ,,obrácená'' teorie her: navrhuje pravidla hry tak, aby racionální hráči svým chováním dosáhli žádoucího výsledku. Příklad jsou \emph{aukce}: anglická (rostoucí, otevřená), holandská (klesající), prvocenová obálková a \emph{Vickreyova} (druhá cena) - u Vickreyovy je dominantní strategií dražit svou pravou hodnotu (pravdomluvnost).
\end{reseni}

\begin{jadro}
Minimax: MAX maximalizuje, MIN minimalizuje hodnoty listu. Alfa-beta ořeže větev při $\alpha\ge\beta$ (až $O(b^{d/2})$). Nashovo ekvilibrium = profil, kde nikdo nezíská jednostrannou změnou (vězňovo dilema: (zrada, zrada)).
\end{jadro}

\kalibrace{Hry (~22\,\%, minimax + teorie her) jsou opakovaná ,,Základy UI'' otázka. Minimax + alfa-beta + Nash = 2-3 body za vysvětlení.}
\past{Alfa-beta nemění \emph{výsledek}, jen rychlost. Nashovo ekvilibrium nemusí být Paretovsky optimální (vězňovo dilema je důkaz).}

\ulohahead{specializace UI-SU}{Bayesovské sítě a pravděpodobnostní inference}
\begin{zadani}
\begin{enumerate}[label=(\alph*)]
  \item \bod{1} Definujte Bayesovskou síť a její vztah k úplné sdružené distribuci (faktorizace).
  \item \bod{2} Jak se síť konstruuje (podmíněná nezávislost) a jak se počítá exaktní inference enumerací?
  \item \bod{3} Popište eliminaci proměnných a aproximační inferenci (Monte Carlo: zamítání, vážení věrohodností).
\end{enumerate}
\end{zadani}

\begin{reseni}
\textbf{(a)} \emph{Bayesovská síť} = orientovaný acyklický graf, kde uzly jsou náhodné proměnné a hrany vyjadřují přímou závislost; každý uzel má tabulku podmíněných pravděpodobností $P(X_i\mid \mathrm{rodiče}(X_i))$. Kóduje \emph{úplnou sdruženou distribuci} faktorizací
\[
P(X_1,\dots,X_n)=\prod_{i=1}^n P\big(X_i\mid \mathrm{rodiče}(X_i)\big),
\]
což šetří parametry oproti plné tabulce (využívá podmíněné nezávislosti).

\textbf{(b)} \emph{Konstrukce}: seřaď proměnné (ideálně příčiny před následky) a pro každou $X_i$ vyber za rodiče $\mathrm{Pa}(X_i)$ takovou podmnožinu předchozích proměnných, aby byla $X_i$ podmíněně nezávislá na ostatních předchozích proměnných při daných $\mathrm{Pa}(X_i)$. \emph{Inference enumerací} dotazu $P(Q\mid e)$: $P(Q\mid e)\propto\sum_{\text{skryté } h}P(Q,e,h)$, kde sčítáme součiny faktoru z faktorizace přes všechny hodnoty skrytých proměnných a nakonec normujeme.

\textbf{(c)} \emph{Eliminace proměnných} sčítá zprava efektivněji: postupně ,,vysčítá'' skryté proměnné, mezivýsledky drží jako faktory a opakované součiny nepřepočítává - výrazně rychlejší než holá enumerace. \emph{Aproximace (Monte Carlo)}: \emph{zamítavé vzorkování} (rejection) generuje vzorky z priorů a zahodí ty neslučitelné s $e$ (neefektivní při vzácném $e$); \emph{vážení věrohodností} (likelihood weighting) fixuje pozorované proměnné a každý vzorek váží součinem jejich pravděpodobností - žádný vzorek nezahazuje.
\end{reseni}

\begin{jadro}
Bayesovská síť = DAG + tabulky $P(X_i\mid\text{rodiče}(X_i))$, kóduje $P(X_1,\dots,X_n)=\prod_i P(X_i\mid\text{rodiče}(X_i))$. Inference: $P(Q\mid e)\propto\sum_{\text{skryté }h}P(Q,e,h)$ (enumerace, efektivněji eliminace proměnných).
\end{jadro}

\kalibrace{Bayesovské sítě (~22\,\% s Bayesovským učením) jsou opakovaná ,,Základy UI'' otázka. Faktorizace + enumerace = 2-bodová záloha; tady je AI rodina, kde je vhodné mít i hlubší zvládnutí.}
\past{Faktorizace je součin $P(X_i\mid\text{rodiče})$ - ne marginálu. Likelihood weighting nezahazuje vzorky (na rozdíl od rejection), proto je efektivnější při vzácné evidenci.}

\ulohahead{specializace UI-SU}{Bayesovské učení a EM algoritmus}
\begin{zadani}
\begin{enumerate}[label=(\alph*)]
  \item \bod{1} Popište Bayesovské učení: prior, věrohodnost, posterior. Co je maximálně věrohodná (ML) a maximálně aposteriorní (MAP) hypotéza a co je Bayesovsky optimální predikce?
  \item \bod{2} Popište EM algoritmus (E-krok a M-krok) pro model se skrytou proměnnou.
  \item \bod{3} Napište vzorec pro responsibility (odpovědnost) v E-kroku a aktualizaci parametru v M-kroku pro skrytou kategoriální proměnnou (shluk).
\end{enumerate}
\end{zadani}

\begin{reseni}
\textbf{(a)} Bayesovo pravidlo: $P(h\mid D)=\frac{P(D\mid h)P(h)}{P(D)}$, kde $P(h)$ je \emph{prior}, $P(D\mid h)$ \emph{věrohodnost}, $P(h\mid D)$ \emph{posterior}. \emph{ML hypotéza} $h_{ML}=\arg\max_h P(D\mid h)$; \emph{MAP} $h_{MAP}=\arg\max_h P(D\mid h)P(h)$ (= ML při uniformním prioru). \emph{Bayesovsky optimální predikce} nepoužívá jedinou hypotézu, ale váží přes všechny: $P(c\mid D)=\sum_h P(c\mid h)P(h\mid D)$.

\textbf{(b)} \emph{EM} maximalizuje věrohodnost při skrytých proměnných střídáním: \emph{E-krok} spočítá očekávané rozdělení skrytých proměnných (responsibility) za daných aktuálních parametru; \emph{M-krok} přepočítá parametry tak, aby maximalizovaly očekávanou (úplnou) log-věrohodnost. Iteruje do konvergence; věrohodnost monotónně neklesá, konverguje do lokálního maxima.

\textbf{(c)} Pro skrytý shluk $z\in\{1,\dots,K\}$ s vahami (priory) $\pi_k$ a modelem $p(x\mid\theta_k)$ je \emph{responsibility} bodu $x_i$ vuči shluku $k$:
\[
r_{ik}=P(z_i=k\mid x_i)=\frac{\pi_k\,p(x_i\mid\theta_k)}{\sum_{j=1}^K \pi_j\,p(x_i\mid\theta_j)}.
\]
\emph{M-krok}: $\pi_k=\frac1n\sum_i r_{ik}$ a parametry $\theta_k$ se odhadnou váženě responsibilitami (např.\ vážený průměr $\mu_k=\frac{\sum_i r_{ik}x_i}{\sum_i r_{ik}}$); u kategoriálních dat = vážené relativní četnosti.
\end{reseni}

\begin{jadro}
$h_{MAP}=\arg\max_h P(D\mid h)P(h)$ ($=$ ML při uniformním prioru). EM: E-krok responsibility $r_{ik}=\frac{\pi_k\,p(x_i\mid\theta_k)}{\sum_j\pi_j\,p(x_i\mid\theta_j)}$; M-krok $\pi_k=\frac1n\sum_i r_{ik}$ a vážené přepočty parametru.
\end{jadro}

\kalibrace{,,Bayesovské učení'' a EM jsou opakovaný ,,Základy UI'' okruh; EM zkouška chtěla \emph{jeden výpočetní update} (responsibility), ne jen pojem. Vzorec responsibility umět zpaměti.}
\past{MAP $\neq$ ML, liší se priorem. EM najde jen \emph{lokální} maximum (závisí na inicializaci); responsibility se musí normovat přes všechny shluky.}

\ulohahead{specializace UI-SU}{Markovské modely a skryté Markovské modely (HMM)}
\begin{zadani}
\begin{enumerate}[label=(\alph*)]
  \item \bod{1} Co je Markovský řetězec a Markovský předpoklad? Definujte skrytý Markovský model (HMM): skryté stavy, přechody, emise.
  \item \bod{2} Popište filtraci, predikci a vyhlazování a uveďte rekurentní vzorec pro filtraci (forward).
  \item \bod{3} Co počítá Viterbiho algoritmus (nejpravděpodobnější průchod)? Jak HMM souvisí s dynamickými Bayesovskými sítěmi?
\end{enumerate}
\end{zadani}

\begin{reseni}
\textbf{(a)} \emph{Markovský řetězec}: posloupnost stavu, kde budoucnost závisí jen na současném stavu (\emph{Markovský předpoklad} $P(X_t\mid X_{1:t-1})=P(X_t\mid X_{t-1})$). \emph{HMM}: skryté stavy $X_t$ s maticí přechodu $P(X_t\mid X_{t-1})$ a pozorování $E_t$ s modelem emisí $P(E_t\mid X_t)$; vidíme jen $E_t$, stavy odhadujeme.

\textbf{(b)} \emph{Filtrace}: $P(X_t\mid e_{1:t})$ (odhad současného stavu z dosavadních pozorování). \emph{Predikce}: $P(X_{t+k}\mid e_{1:t})$ (budoucí stav). \emph{Vyhlazování}: $P(X_k\mid e_{1:t})$ pro $k<t$ (minulý stav s využitím pozdějších pozorování). Rekurence filtrace (forward):
\[
P(X_t\mid e_{1:t})\ \propto\ P(e_t\mid X_t)\sum_{x_{t-1}}P(X_t\mid x_{t-1})\,P(x_{t-1}\mid e_{1:t-1}),
\]
tj.\ krok predikce (přes přechod) a poté přenásobení emisí a normalizace.

\textbf{(c)} \emph{Viterbi} dynamickým programováním najde \emph{nejpravděpodobnější posloupnost skrytých stavu} pro dané pozorování (max místo součtu ve forward rekurzi + zpětný chod). HMM je speciální případ \emph{dynamické Bayesovské sítě} (řetězec s jednou skrytou proměnnou na krok); DBN dovolují víc proměnných a strukturovanějších závislostí na časový krok.
\end{reseni}

\kalibrace{HMM (~12\,\%) je opakované ,,Základy UI'' téma (filtrace se zkoušela numericky). Definice + forward rekurence = 2-bodová záloha; pro 3 body umět jeden krok filtrace dosadit.}
\past{Filtrace je jen ,,dopředu'', vyhlazování využívá i \emph{budoucí} pozorování. Viterbi hledá nejpravděpodobnější \emph{celou cestu}, ne stav po stavu.}

\ulohahead{specializace UI-SU}{Markovské rozhodovací procesy (MDP) a zpětnovazební učení}
\begin{zadani}
\begin{enumerate}[label=(\alph*)]
  \item \bod{1} Definujte MDP (stavy, akce, přechody, odměny, diskont) a pojmy strategie (policy) a užitek (hodnota) stavu.
  \item \bod{2} Napište Bellmanovu rovnici pro optimální hodnotu a popište iteraci hodnot (value iteration).
  \item \bod{3} Co je iterace strategií? Popište zpětnovazební učení: Q-učení (update) a kompromis explorace vs exploitace.
\end{enumerate}
\end{zadani}

\begin{reseni}
\textbf{(a)} \emph{MDP} $=(S,A,P,R,\gamma)$: stavy $S$, akce $A$, přechody $P(s'\mid s,a)$, odměny $R(s,a,s')$, diskont $\gamma\in[0,1)$. \emph{Strategie} $\pi:S\to A$ určuje akci v každém stavu. \emph{Hodnota} (užitek) stavu při $\pi$: očekávaná diskontovaná suma odměn $V^\pi(s)=\mathbb{E}[\sum_{t\ge0}\gamma^t R_t\mid s_0=s,\pi]$.

\textbf{(b)} \emph{Bellmanova rovnice} pro optimální hodnotu:
\[
V^*(s)=\max_{a}\sum_{s'}P(s'\mid s,a)\big[R(s,a,s')+\gamma V^*(s')\big].
\]
\emph{Iterace hodnot}: inicializuj $V$, opakovaně aplikuj pravou stranu jako přiřazení (Bellmanuv update) pro všechny stavy, dokud se $V$ nestabilizuje; pak optimální strategie $\pi^*(s)=\arg\max_a\sum_{s'}P[\,R+\gamma V^*\,]$.

\textbf{(c)} \emph{Iterace strategií}: střídá \emph{vyhodnocení} dané strategie (vyřeš $V^\pi$) a \emph{zlepšení} (greedy vuči $V^\pi$); konverguje v konečně krocích. \emph{Zpětnovazební učení} (neznámé $P,R$): \emph{Q-učení} aktualizuje akční hodnotu z pozorovaného přechodu
\[
Q(s,a)\leftarrow Q(s,a)+\alpha\big[r+\gamma\max_{a'}Q(s',a')-Q(s,a)\big]
\]
(off-policy). \emph{Explorace vs exploitace}: je třeba občas zkusit i nepreferované akce (např.\ $\varepsilon$-greedy), jinak agent uvázne v suboptimu; SARSA je on-policy obdoba (místo $\max$ použije skutečně zvolenou $a'$).
\end{reseni}

\kalibrace{MDP/RL (~15\,\%) je opakovaná ,,Základy UI'' otázka (Bellmanova rovnice se zkoušela). Definice MDP + Bellman + iterace hodnot = 2-bodová záloha.}
\past{Bellman pro optimum má $\max_a$ (u dané strategie je tam suma přes $\pi$). Q-učení je off-policy ($\max$), SARSA on-policy - nezaměňuj.}

\ulohahead{specializace UI-SU}{Logické uvažování: DPLL, řetězení, rezoluce}
\begin{zadani}
\begin{enumerate}[label=(\alph*)]
  \item \bod{1} Zopakujte CNF a popište algoritmus DPLL pro řešení splnitelnosti (SAT): jednotková propagace a čistý literál.
  \item \bod{2} Co jsou Hornovy klauzule a jak funguje dopředné a zpětné řetězení?
  \item \bod{3} Popište rezoluci jako rozhodovací/dokazovací metodu (důkaz sporem, prázdná klauzule). K čemu je v predikátové logice unifikace?
\end{enumerate}
\end{zadani}

\begin{reseni}
\textbf{(a)} \emph{CNF} = konjunkce klauzulí (disjunkcí literálu). \emph{DPLL} prohledává přiřazení s prořezáváním: (1) \emph{jednotková propagace} - klauzule s jediným literálem vynutí jeho hodnotu (a zjednoduší ostatní); (2) \emph{čistý literál} - proměnná vyskytující se jen v jedné polaritě se nastaví tak, aby tyto klauzule splnila; (3) jinak rozděl podle nějaké proměnné (větev true/false) a rekurzivně. Konec: prázdná množina klauzulí = splnitelné, prázdná klauzule = konflikt (backtrack).

\textbf{(b)} \emph{Hornova klauzule} má nejvýše jeden pozitivní literál (tj.\ tvar $a_1\wedge\dots\wedge a_k\Rightarrow b$). \emph{Dopředné řetězení}: z známých faktu opakovaně odvozuj hlavy pravidel, jejichž těla jsou splněna, dokud nepřibývají nové fakty (data-driven). \emph{Zpětné řetězení}: od dotazovaného cíle rekurzivně dokazuj těla pravidel, která ho mají v hlavě (goal-driven, základ Prologu). Pro Hornovy klauzule je to lineární/efektivní.

\textbf{(c)} \emph{Rezoluce}: dokazuje $T\models\varphi$ sporem - k $T$ přidá $\neg\varphi$, převede na CNF a opakovaně tvoří rezolventy; odvození \emph{prázdné klauzule} znamená nesplnitelnost, tedy platnost $\varphi$. V \emph{predikátové} logice je nutná \emph{unifikace}: najití nejobecnějšího substitučního přiřazení proměnných, které dva literály ztotožní, aby šly rezolvovat.
\end{reseni}

\begin{jadro}
DPLL = jednotková propagace + čistý literál + větvení s backtrackingem. Hornovy klauzule ($\le1$ pozitivní literál) = dopředné/zpětné řetězení. Rezoluce = důkaz sporem (přidej $\neg\varphi$, odvoď prázdnou klauzuli); v predikátové logice s unifikací.
\end{jadro}

\kalibrace{Logické uvažování (~10\,\%, DPLL se zkoušelo) je vzácnější ,,Základy UI'' okruh. CNF + DPLL kroky = 2-bodová záloha; navíc se prolíná s matematickou logikou.}
\past{Jednotková propagace a čistý literál nejsou totéž. Rezoluce dokazuje sporem (přidej negaci cíle); cílem je odvodit prázdnou klauzuli.}

\ulohahead{specializace UI-SU}{Automatické plánování a reprezentace znalostí}
\begin{zadani}
\begin{enumerate}[label=(\alph*)]
  \item \bod{1} Popište formulaci plánovacího problému a definici operátoru (předpoklady a efekty), počáteční stav a cíl.
  \item \bod{2} Vysvětlete dopředné (progression) a zpětné (regression) plánování a jejich rozdíl.
  \item \bod{3} Co je situační kalkulus a problém rámce (frame problem) a jak se řeší?
\end{enumerate}
\end{zadani}

\begin{reseni}
\textbf{(a)} \emph{Plánovací problém} (např.\ STRIPS): počáteční stav (množina platných faktu), množina \emph{operátoru} a cílová podmínka. \emph{Operátor} má \emph{předpoklady} (precondition - co musí platit, aby šel provést) a \emph{efekty} (add list = co začne platit, delete list = co přestane platit). Plán = posloupnost operátoru vedoucí z počátečního stavu do stavu splňujícího cíl.

\textbf{(b)} \emph{Dopředné plánování}: začni v počátečním stavu a aplikuj použitelné operátory, dokud nedosáhneš cíle (prohledávání stavového prostoru vpřed); velký větvící faktor. \emph{Zpětné plánování}: začni od cíle a hledej operátory, jejichž efekty cíl (částečně) splňují, a nahraď cíl jejich předpoklady (regrese); pracuje jen s relevantními operátory, ale stavy jsou částečné (množiny podcílu).

\textbf{(c)} \emph{Situační kalkulus} popisuje svět v predikátové logice: \emph{fluenty} (vlastnosti závislé na situaci, např.\ $At(x,s)$), akce a funkci $do(a,s)$ pro situaci po akci. \emph{Problém rámce}: je třeba vyjádřit nejen co se akcí \emph{změní}, ale i co zustane \emph{beze změny} - jinak by se počet axiomu rozrostl pro každou dvojici akce/fluent. Řeší se \emph{successor-state axiomy} (pro každý fluent jeden axiom popisující, kdy v nové situaci platí), což stručně zachytí setrvačnost (co akce neovlivní, zustává).
\end{reseni}

\begin{jadro}
Operátor $=$ (precondition, add efekty, delete efekty). \emph{Dopředné} plánování jde od počátku vpřed, \emph{zpětné} regreduje od cíle přes efekty operátoru. Problém rámce $=$ jak vyjádřit, co se akcí nemění (successor-state axiomy).
\end{jadro}

\kalibrace{Plánování a reprezentace znalostí jsou vzácnější ,,Základy UI'' rodiny, ale legální draw do skoro vždy přítomného AI slotu. Bez této zálohy hrozí na takové otázce skóre 0-1, což přímo ohrožuje podlahu specializace ($\ge7/12$).}
\past{Dopředné vs zpětné nepleť: regrese jde od cíle a nahrazuje cíl předpoklady operátoru. Problém rámce není o ,,výpočtu'', ale o stručném popisu neměnnosti.}


\end{document}
